\documentclass[11pt,openany]{book}

\usepackage{iftex}
\ifPDFTeX
  \usepackage[T1]{fontenc}
  \usepackage[utf8]{inputenc}
  \usepackage{lmodern}
\else
  \usepackage{lmodern}
  \usepackage{fontspec}
\fi
\usepackage[margin=1.15in]{geometry}
\usepackage{microtype}
\usepackage{amsmath,amssymb,amsthm}
\usepackage{booktabs}
\usepackage{tabularx}
\newcolumntype{Y}{>{\raggedright\arraybackslash}X}
\usepackage{enumitem}
\usepackage{tikz}
\usepackage[hidelinks]{hyperref}
\setlength{\emergencystretch}{2em}

\theoremstyle{definition}
\newtheorem{definition}{Definition}[chapter]
\newtheorem{hypothesis}{Hypothesis}
\newtheorem{protocol}{Protocol}[chapter]

\newcommand{\todo}[1]{\par\noindent\textit{[Draft note: #1]}\par}

\title{\textbf{Phase Fraction and Filament Geometry\\in Structured Foods}\\[1em]
\large A Configurable Process Architecture for Emulsions,\\Core--Shell Strands, and Consolidated Textures}
\author{Flyxion\\Independent Researcher}
\date{September 2026}

\begin{document}

\frontmatter
\maketitle

\chapter*{Abstract}
\addcontentsline{toc}{chapter}{Abstract}

This monograph proposes a configurable food manufacturing platform in which characterized ingredient streams are routed through a shared set of transformation and forming operations to produce several classes of food. Its principal contribution is a forming method: composite strands built from a soft core and a formable outer layer, drawn to a target diameter, coiled together, and pressed into a larger structure. Among the variables that shape texture, two structural ones are given particular attention because they can be adjusted independently of formulation. At the strand scale, the core fraction of a composite filament, together with strand diameter, coiling geometry, and the degree of bonding between strands, is proposed to move a food between noodle-like, shredding, and dense sliceable forms. At the droplet scale, the dispersed phase fraction of an emulsion offers an analogy, not a recipe continuum: milk, cream, and mayonnaise differ in more than their fat content, but the comparison isolates the idea that the proportion of an inner phase can change bulk behavior sharply. Layered and fibrous foods occupy one region of the resulting design space; reconstructed fruits, dairy-like products, breads, and confections occupy others. Animal-derived foods appear as reference structures where comparison is needed, not as targets to be imitated. The work separates practices that have been carried out domestically from methods that remain proposals, states its central claims as testable hypotheses, and specifies a first experimental sequence.

\chapter*{Status of the Claims}
\addcontentsline{toc}{chapter}{Status of the Claims}

The material in this book falls into three categories, and each chapter marks which one it draws on.

\begin{description}[leftmargin=1.5em]
\item[Practiced.] Domestic methods repeated over several years: hydrating textured vegetable protein, lentils, or bread and frying them in broth or oil with onion, herbs, and varied spice mixtures to change the perceived identity of a dish.
\item[Proposed.] Methods not yet built or tested: centrifugal characterization pipelines, core--shell strand formation, draw-down, coiling, press consolidation, rolling hotplate compression, and reconstructed produce.
\item[Hypothesized.] Claims about what the proposed methods will do, stated below and tested in Part VIII.
\end{description}

\begin{hypothesis}[Late binding]\label{h:late}
Beyond a threshold of flavor complexity, the perceived identity of a savory dish is determined mainly by its final aroma, sauce, and fat profile, and tasters cannot reliably identify the protein or starch base beneath it.
\end{hypothesis}

\begin{hypothesis}[Phase fraction as a control variable]\label{h:phase}
In two-phase food structures, the fraction of the inner phase is an independently adjustable structural variable whose effect on cross-sectional structure and directional mechanical response can be measured while formulation, setting, orientation, and bonding are held fixed.
\end{hypothesis}

\begin{hypothesis}[Draw invariance]\label{h:draw}
A core--shell preform can be drawn to a much smaller diameter while approximately preserving its core fraction.
\end{hypothesis}

\begin{hypothesis}[Shell fusion]\label{h:fusion}
Under suitable pressure and temperature, the shells of adjacent coiled strands fuse into a connected network while the cores remain distinct, and this internal structure survives cooking.
\end{hypothesis}

None of these is established by the domestic practice. Hypothesis~\ref{h:late} is suggested by it; Hypotheses~\ref{h:phase} through~\ref{h:fusion} are engineering conjectures.

The first structural prediction of the book is correspondingly limited. The proposed process should be capable of producing strands with measurable core and shell regions that can be formed, coiled, and consolidated into blocks; the experiments test whether changing core fraction and coiling geometry produces measurable changes in the cross-sectional structure and directional mechanical response of those blocks. Sectioning before and after cooking, recording leakage and breakage, and comparing shear force along and across the coiling direction test this prediction directly. Perceptual identity and nutritional equivalence are separate questions, taken up in their own chapters and not presupposed by the structural work.

\chapter*{Notation}
\addcontentsline{toc}{chapter}{Notation}

\begin{tabularx}{\linewidth}{@{}lX@{}}
\toprule
Symbol & Meaning \\
\midrule
$\psi$ & Dispersed phase volume fraction of an emulsion \\
$\phi$ & Core area fraction of a composite strand, $(r_c/R)^2$ \\
$R,\ r_c$ & Strand outer radius, core radius \\
$d$ & Final strand diameter \\
$\lambda$ & Draw ratio, preform diameter divided by final diameter \\
$\theta$ & Coiling helix angle measured from the mandrel axis \\
$n$ & Number of coiled layers \\
$P,\ T,\ t$ & Press pressure, temperature, dwell time \\
$C,\ S$ & Core and shell formulations \\
$F$ & Finishing profile (surface and interlayer seasoning, oil, crust) \\
$\mathcal{A}$ & Anisotropy index, shear force along the grain divided by shear force across it \\
$f$ & Fusion index, peak separation force of a pressed strand pair divided by peak tensile break force of a single strand, under a stated geometry (dimensionless, operational) \\
$\eta$ & Packing fraction of strands before consolidation \\
\bottomrule
\end{tabularx}

\tableofcontents

\mainmatter

%======================================================================
\part{Two-Phase Food Architecture}
%======================================================================

\chapter{Carrier, Matrix, Inclusion}\label{ch:carrier}

Status: practiced as a way of cooking; proposed as a descriptive framework.

\section{The fried starch carrier}

A large share of the world's prepared dishes can be described with three terms. A \emph{carrier} is a starch-based body that supplies bulk and bite: bread, rice, noodles, flatbread, dumpling wrappers, potato. A \emph{matrix} is the sauce, gravy, or binding phase that coats or fills the carrier. \emph{Inclusions} are the discrete pieces distributed through the matrix: vegetables, legumes, cheese, meat, herbs.

Pizza is a baked carrier with a tomato matrix and cheese inclusions. Curry is a spiced matrix served over a rice or flatbread carrier. Spaghetti is a boiled carrier coated in a matrix. A sandwich, a burrito, a dumpling, and a pie differ mainly in how carrier and matrix are arranged relative to each other: stacked, wrapped, enclosed, or layered.

The carrier is very often fried or baked, and in either case its surface is browned. Frying and baking do the same work at the surface: they drive off water, raise the temperature well above boiling, and produce crust and roasted flavor. Seen this way, fried bread is not a special category but the common ancestor of many dishes, and much of what distinguishes those dishes lies in the matrix and the inclusions rather than in the carrier.

\section{Sauce and matrix}

The matrix carries most of the salt, acid, fat, and aroma of a dish. It is also the easiest component to change. A tomato matrix becomes a curry matrix with a different spice mixture and fat; a gravy becomes a cheese sauce with a different thickener and flavor base. Because the matrix is liquid or semi-liquid when prepared, it can be adjusted continuously by tasting, which is why cooks learn to correct a dish at the sauce stage.

\section{Binders shared across dishes}

Some foods that appear to be made mainly of meat are in fact partly carrier. Meatloaf and many sausages contain bread, breadcrumbs, or other starch that binds the mixture and holds moisture. Donair and gyro meat are formed from ground meat that is often combined with binders and spices and compressed into a cone. Potato, oats, and rice serve the same role in other traditions.

These binders show that the boundary between carrier and inclusion is not fixed. A food recognized as meat can contain a substantial proportion of starch and still be recognized as meat, because recognition depends on its matrix, fat, aroma, and texture as much as on its protein source. Chapter~\ref{ch:config} states this as a hypothesis.

\section{Transitional dishes}

If dishes are combinations of carrier, matrix, and inclusions, then new dishes can be made by moving one component at a time. A curry matrix on a pizza carrier, a sweet dairy matrix on a savory spiced base, a chili thickened and set like a custard: each is a step from a known dish toward another. Chapter~\ref{ch:transitional} treats this as a design space and asks which steps hold together and which fail.\footnote{This framing is related to an older satirical treatment of cooking as a hub of interchangeable parts. It is used here as a working description, not as a claim that every cuisine reduces to three terms.}

\section{Relevance to the rest of the book}

The framework matters for manufacturing because it suggests where shared equipment can serve many products. If carriers are a small number of starch forms, matrices are adjustable liquids, and inclusions are formed pieces, then a platform that can make a few carrier types, mix and adjust matrices, and form structured inclusions can make a very large number of dishes. The composite-strand method of Part IV is one way of making structured inclusions and carriers. It is not the only thing such a platform does.

%----------------------------------------------------------------------
\chapter{Perceptual Coordinates}\label{ch:coords}

Status: summarizes established sensory science.

\section{Texture coordinates}

Texture is perceived through the teeth, tongue, and jaw muscles as a food is bitten and chewed. Sensory and instrumental work usually describes it with a small set of properties: hardness (force to compress), cohesiveness (how well the food holds together), springiness (recovery after compression), chewiness (work to reduce it for swallowing), and fracture behavior (whether it snaps, crumbles, or tears)~\cite{bourne2002}. Direction also matters. A food with grain resists differently along and across that grain, and anisotropy is itself a perceived property.

\section{Basic tastes}

The tongue detects a small number of basic tastes: sweet, salty, sour, bitter, and umami, the savory taste associated with glutamate and certain nucleotides. These are coarse channels. They tell a taster that a food is salty or sour but carry little information about what the food is.

\section{Fat as carrier and coordinate}

Fat contributes to perception in three ways. It has its own mouthfeel: coating, smoothness, lubrication. It melts, so its behavior changes with temperature in the mouth. And it dissolves many aroma compounds, holding them in the food and releasing them during chewing. Fat therefore acts both as a coordinate in its own right and as the vehicle for a large share of aroma.

\section{Aroma and retronasal perception}

Aroma reaches the nose by two routes. Sniffing draws volatile compounds through the nostrils. Chewing and swallowing push them from the mouth up the back of the throat into the nasal cavity, a route called retronasal. Retronasal aroma is experienced as taste, not smell, which is why food seems tasteless when the nose is blocked. The distinction between the two routes, and the tendency to experience retronasal smell as taste, was set out in an influential paper by Rozin~\cite{rozin1982}. Most of what people describe as the flavor of a food, as opposed to its basic tastes, is aroma.

Aroma is also the most informative channel. There are only a few basic tastes but a very large number of volatile compounds, and their combinations are what allow a taster to tell beef from chicken, or cumin from coriander.\footnote{This is why the list of coordinates in this chapter treats aroma as a family of axes rather than a single one. It cannot be summarized by one number.}

\section{Coordinates for design}

For the purposes of this book, a food is described by texture coordinates, basic taste levels, fat content and state, and an aroma profile. The list is not a complete account of perception. Temperature, appearance, sound, and expectation all matter. It is a working set of variables that a process can control and that a sensory panel can report on.

%----------------------------------------------------------------------
\chapter{Configural Perception and Identity}\label{ch:config}

Status: draws on established research; states Hypothesis~\ref{h:late}.

\section{Mixture perception limits}

Odors in mixtures are often not perceived as lists of their components. In a study of odor mixture identification, Laing and Francis found that people, including subjects trained on the individual odors, had great difficulty identifying the components of mixtures, and rarely identified more than three or four, even though each odor was familiar by itself~\cite{laing1989}. A separate line of research describes how mixtures are perceived rather than how well they can be analyzed. A review by Thomas-Danguin and colleagues describes a range between elemental perception, in which at least some components remain perceptible within the mixture, and configural perception, in which the mixture is perceived as a single new odor distinct from its components; it reports that more complex mixtures are more inclined to be perceived in the configural way~\cite{thomasdanguin2014}.

Their mixtures were small sets of single odorants presented under laboratory conditions. The experiment supports the general limit on analyzing mixtures. It does not show that a dish seasoned with ten to fifty spices conceals the protein or starch beneath it. That extension is a separate claim, and it remains Hypothesis~\ref{h:late}, to be tested in Experiment E1.

Food aromas are far more complex than these laboratory mixtures, often involving dozens or hundreds of volatile compounds. It is plausible, though not established by these studies, that what is recognized in a food aroma is its configural character rather than a set of ingredients.

\section{Complexity thresholds in cooking}

A spice mixture of ten to fifty components, as commonly used in curries and spice blends, is well beyond the identification limit. A cook adjusting such a mixture changes its overall character, not a set of independently perceived notes. This explains a familiar experience: a dish can be recognized as a particular curry, chili, or stew without anyone being able to say which ingredients produce its character.

It also suggests that when enough is happening in the matrix, differences in the underlying carrier or protein may be hard to perceive. Fried textured vegetable protein, lentils, and bread, each cooked in oil with onion, herbs, and a complex spice mixture, may be perceived primarily as the dish they are served in.

\section{The hypothesis}

Hypothesis~\ref{h:late} states this possibility in testable form: beyond a threshold of flavor complexity, perceived identity is determined mainly by final aroma, sauce, and fat, and tasters cannot reliably identify the base beneath it.

The hypothesis is suggested by domestic practice, but that practice does not establish it. Cooks and the people they cook for know what went into a dish, and expectation shapes perception. The test in Experiment E1 removes that knowledge: bases are prepared identically, coded, and presented in triangle tests, in which a taster must pick the odd sample out of three.

\section{Where the hypothesis would fail}

The hypothesis predicts its own failure conditions, and they are informative.

\begin{itemize}
\item \textbf{Texture differences.} If the bases differ strongly in chew or fracture, tasters may distinguish them by texture even when aroma is identical. This would show that texture coordinates carry identity independently of aroma, which is an argument for the strand method rather than against it.
\item \textbf{Low complexity.} With simple seasoning, the base's own flavor should be detectable. Varying the number of spice components across trials would locate the threshold.
\item \textbf{Characteristic off-notes.} Some bases carry distinctive flavors, such as the beany notes of some legumes, that may survive seasoning.
\end{itemize}

A result in which tasters can tell bases apart only under some of these conditions would map where the late-binding claim holds, which is more useful than a simple confirmation.

%----------------------------------------------------------------------
\chapter{Meatiness as a Reaction Class}\label{ch:meaty}

Status: summarizes established flavor chemistry.

\section{Browning chemistry}

When sugars and amino acids are heated together, they undergo the Maillard reaction, a large network of reactions that produces brown pigments and hundreds of volatile compounds. It is responsible for the flavor of bread crust, roasted coffee, toasted nuts, fried onion, and seared meat. The products depend on which sugars and amino acids are present, on temperature, on water content, and on time.

\section{Sulfur and the savory note}

The aromas most associated with cooked meat come largely from sulfur-containing compounds formed when sulfur-bearing amino acids, especially cysteine, react with sugars during heating. Some of these compounds are perceived as meaty at extremely low concentrations. The same reactions can be run without meat: heating cysteine with a reducing sugar such as ribose produces a strong meat-like aroma, and so-called reaction flavors made in this way are used in commercial stocks, gravies, and seasonings. The chemistry of meat flavor formation, including the roles of Maillard reactions, sulfur-containing precursors, and lipid degradation, is reviewed by Mottram~\cite{mottram1998}.

Onions and garlic contain sulfur compounds of their own, and frying them adds sulfur-derived and browned notes together. This is one reason why frying almost any starch or protein with onion in oil produces an aroma that people associate with meat cooking.

\section{Roast bitterness}

Roasting and searing also produce bitter and toasty compounds, including many shared by coffee, cocoa, and the crust of seared meat. A small amount of coffee or cocoa in a savory dish supplies some of that roasted character without searing. The use of cocoa in some chili and mole traditions is an established example.

\section{Heme as catalyst}

Heme, the iron-containing molecule that gives blood and red meat their color, has been promoted as the key to meat flavor in some plant-based products. Its role is better described as catalytic: iron speeds up reactions during cooking that form meat-like aromas. It is one route into the reaction class, not the source of meatiness itself.\footnote{The claim here is not that heme has no effect, but that meaty aroma can be reached by other routes, as the existence of heme-free reaction flavors shows.}

\section{Consequence for process design}

If meatiness is a class of cooking reactions, it belongs to the finishing stage of the process, where sugars, sulfur-bearing amino acids, fat, and heat are brought together at the surface of a formed food. It does not have to be built into the structure. This supports the division of labor proposed in Part IV: structure sets texture, and finishing sets aroma.

%----------------------------------------------------------------------
\chapter{Emulsions and the Dispersed Phase}\label{ch:emulsion}

Status: summarizes established colloid science; the comparison with strands is an analogy.

\section{Oil-in-water structure}

An emulsion is a dispersion of droplets of one liquid in another with which it does not mix. In an oil-in-water emulsion, oil droplets are suspended in a continuous water phase. The droplets tend to merge and separate from the water, and they are kept apart by molecules at their surfaces, typically proteins or phospholipids, that form a protective film.

Milk, cream, mayonnaise, many sauces, and many spreads are oil-in-water emulsions. The fraction of the volume occupied by the dispersed droplets, written here $\psi$, is one of the variables that differs among them. McClements gives a comprehensive treatment of food emulsions, including their formation, stability, and rheology~\cite{mcclements2015}.

\section{Milk and mayonnaise}

Whole milk contains a few percent fat by weight in fine droplets, and it pours. Cream contains much more fat and is thicker. Mayonnaise is typically more than two-thirds oil by weight, held as droplets in a small amount of water phase stabilized by egg yolk, and it holds its shape on a spoon. Arranged by fat content, the three suggest a scale along which a liquid becomes a spread as $\psi$ increases.

The comparison is an analogy, not a recipe continuum. The three foods differ in more than droplet fraction. Milk's structure also depends on casein micelles, whey proteins, and lactose; mayonnaise depends on egg yolk emulsifiers, acid, and salt; droplet sizes differ; and the processes that make them are different. Adding oil to milk does not make mayonnaise. The analogy is useful because it isolates one structural idea: that the proportion of an inner phase can change bulk behavior dramatically.

\section{Packing and the onset of solid behavior}

The change from pourable to shape-holding has a geometric explanation. When droplets occupy a small fraction of the volume, they move freely past one another. As $\psi$ rises, droplets crowd together, and above a certain packing they can no longer move without deforming. For spheres of equal size in random arrangement, this crowding becomes severe at roughly two-thirds of the volume. Beyond it, droplets are pressed into polyhedral shapes separated by thin films, and the emulsion resists flow and holds its shape. Mason, Bibette, and Weitz measured this transition directly, showing that concentrated emulsions develop an elastic modulus that rises steeply once droplets are compressed beyond close packing~\cite{mason1995}.\footnote{Real emulsions have a range of droplet sizes, which allows denser packing before deformation. The two-thirds figure is a reference value, not a threshold for any particular food.}

The same geometry reappears in Chapter~\ref{ch:consolidation}, where round strands pressed together deform into polygons that fill space. The resemblance is structural: in both cases a dispersed or enclosed phase packed beyond a certain fraction forces the continuous phase into thin shared walls.

\section{Emulsions as strand cores}

Emulsions have a direct role in the strand process as core materials. A high-$\psi$ emulsion is soft, holds its shape, and carries fat and fat-soluble aroma. As a core it would supply juiciness and flavor to the center of a strand block. Its behavior during drawing and cooking, especially whether it remains emulsified or separates into free oil, is a question for the experiments.

\section{The phase-fraction hypothesis}

Hypothesis~\ref{h:phase} is deliberately limited. It does not claim that inner-phase fraction is the primary determinant of texture, since formulation, interfacial behavior, setting, orientation, and bonding may all matter as much or more. It claims that inner-phase fraction is an independently adjustable variable whose effect can be measured while the others are held fixed. In emulsions, varying $\psi$ with a fixed emulsifier and droplet size shows that effect clearly. In strands, varying $\phi$ with fixed core and shell formulations and fixed setting, coiling, and pressing conditions is the corresponding test.

%----------------------------------------------------------------------
\chapter{Transitional Foods and Unstable Regions}\label{ch:transitional}

Status: proposed as a research program.

\section{Interpolation between dishes}

If dishes are described by their carrier, matrix, inclusions, and perceptual coordinates, then any two dishes are points in a common space and there are paths between them. Chili and cheesecake are far apart: one is a savory, spiced, chunky stew, the other a sweet, smooth, set dairy custard on a crumb base. But there are dishes between them. A chili set with egg and dairy like a savory custard, a cheesecake with warm spices and a savory crust, a sweet bean dessert: each is a step along a path.

Many established foods sit at such intermediate points. Sweet and savory combinations, bean desserts, spiced cakes, and savory custards are examples from many cuisines. They show that the space between familiar dishes is not empty.

\section{Failure modes}

Not every intermediate point is a stable food. Moving between dishes can cross into regions where the combination breaks down physically or perceptually.

\begin{description}[leftmargin=1.5em]
\item[Curdling.] Acid and heat cause dairy proteins to clump and separate. A tomato- or vinegar-based matrix combined with milk or cream at the wrong temperature curdles.
\item[Fat separation.] An emulsion pushed beyond its stabilizer's capacity, or heated too far, breaks and releases free oil.
\item[Texture conflict.] Some combinations of textures read as defects: a grainy component in a smooth one, a soggy carrier under a wet matrix.
\item[Perceptual conflict.] Some aroma and taste combinations are judged unpleasant regardless of physical stability, and these judgments vary among people and cultures.
\end{description}

\section{Mapping the stable regions}

These failure modes make the space of foods uneven: there are regions of stable, acceptable combinations and regions where combinations fail. Mapping them is an experimental program. A path between two dishes is divided into steps; at each step a sample is prepared, its physical stability recorded (separation, curdling, texture), and its acceptability rated. Where failures appear, the formulation is adjusted, for example by changing the order of addition, the stabilizer, or the temperature, to see whether the failure can be moved or removed.\footnote{Stabilizers such as starches, gums, and emulsifiers exist largely to extend the stable region of a formulation. Part of the program is identifying which failures they can repair and which they cannot.}

\section{Relevance to the platform}

The platform described in Part III makes a large number of combinations possible. Knowing which regions of the space are stable determines which of those combinations are worth making. The transitional-food program is therefore not a curiosity but a guide to product design: it identifies where a shared set of operations can produce acceptable food and where it cannot.

%======================================================================
\part{Precedents}
%======================================================================

\chapter{Industrial Fractionation}\label{ch:industrial}

Status: summarizes established industrial practice.

Taking agricultural materials apart and recombining the pieces is not a new idea. It is the basis of several of the largest food industries. This chapter reviews four of them, both to show that the unit operations proposed later already work at scale and to state precisely what this book does and does not claim as new.

\section{Wet milling}

Corn wet milling separates kernels into their main components. Kernels are steeped in warm water, softening them and loosening their structure. They are then coarsely ground, and the germ, which is rich in oil, is separated. Further grinding and screening remove the fiber of the hull, and centrifugal separation divides what remains into a protein-rich fraction and a starch slurry. Each stream becomes a product or an input: oil from the germ, feed from the fiber and protein, and starch that is sold as such or converted into sweeteners and fermentation feedstocks. Laboratory and pilot-scale wet-milling procedures, and the variables that govern component yields, are reviewed by Singh and Eckhoff~\cite{singh1996}.

Wet milling shows three things relevant here. A single crop can feed many product lines. Centrifugal separation of a wet slurry by density is a mature, continuous operation. And every fraction has to go somewhere, so the economics of the whole depend on finding uses for all of them.

\section{Dairy cracking and recombination}

Milk is routinely separated and reassembled. Centrifugal cream separators divide whole milk into cream and skim milk. Skim milk can be acidified or treated with rennet to separate casein curd from whey. Whey can be filtered through membranes to concentrate its proteins and separate lactose. Casein, whey proteins, lactose, and milk fat are then sold as ingredients and recombined in other foods.

The dairy industry also reverses the process. Standardized milk is made by separating whole milk centrifugally into cream and skim milk and then blending cream back into the skim in controlled proportions, so that the product has a specified fat content~\cite{tetrapak}. Recombined milk goes further. It is made by adding water to skim milk powder and adding milk fat separately in the quantity needed for the desired fat content, then emulsifying and homogenizing the mixture, producing standardized liquid milk from stored dairy components~\cite{tetrapak}.

The dairy industry's own vocabulary marks the distinction that matters here. Reconstituted milk is milk powder with water added back. Recombined milk is built from separate components: the skim powder supplies the non-fat solids and the fat is supplied independently~\cite{tetrapak}. Recombination is therefore an established industrial example of the approach this book generalizes: separating a food into components that carry distinct functions, storing them, and rebuilding a product to specification.

\section{Soy processing}

Soybeans are dehulled, flaked, and extracted to remove oil, leaving defatted flakes. These can be ground into flour, washed to remove soluble sugars and produce protein concentrate, or dissolved and precipitated to produce protein isolate. Textured vegetable protein is made by extruding defatted soy flour under heat and pressure so that it expands into a porous, chewy structure. Soy milk and tofu production leave a fibrous residue, okara, which has uses of its own.

Textured vegetable protein matters to this book in a specific way. It is a fractionated and restructured product that is cheap, shelf-stable, and widely used in domestic cooking. It is the base used in the practiced methods described in Part I.

\section{Biorefineries}

The biorefinery concept extends the logic of wet milling to plant material generally: a crop or residue is separated into as many valuable fractions as possible, by analogy with the way a petroleum refinery produces many products from crude oil. The concept is well developed in the literature for fuels, chemicals, and materials, and less so for food at small or regional scale.\footnote{Much biorefinery work treats food and feed as low-value outlets for fractions left after higher-value chemicals are extracted. The priority is reversed here: food is the main product.}

\section{What is and is not new here}

None of the separation operations described in Part III is new, and none of the forming operations in isolation is new. Extrusion, coating, drawing, coiling, and pressing are each used in food or other industries.

The proposals of this book lie elsewhere:

\begin{enumerate}
\item a shared, configurable set of operations serving several product classes, rather than a dedicated line per commodity;
\item operation at small or regional scale, with locally available crops treated as interchangeable inputs;
\item the combination of coaxial preform construction, drawing, coiling, and consolidation into a method for building layered and fibrous foods with independently adjustable structural variables;
\item batch records and release decisions applied uniformly across the whole platform.
\end{enumerate}

Of these, the third is the most specific and the one on which the experimental program concentrates.

%----------------------------------------------------------------------
\chapter{Existing Structuring Methods}\label{ch:structuring}

Status: summarizes established methods.

Several methods already produce fibrous or layered structure from plant proteins. Each controls some variables and leaves others to chance. The strand method is proposed to fill a gap among them, and that gap is easiest to see by comparison.

\section{Extrusion and cooling dies}

Low-moisture extrusion forces a relatively dry protein mixture through a hot die, where the sudden release of pressure makes it expand into a porous product. This is how textured vegetable protein is made. The resulting texture is spongy and becomes chewy when rehydrated, but its internal organization is largely random.

High-moisture extrusion works with wetter mixtures and passes the hot, plasticized material through a long cooling die. As the material flows and cools in the die, its proteins align and set, producing a fibrous, layered structure. High-moisture extrusion is widely used to make whole-cut plant-based products. How its fibrous structure forms is still debated; van der Sman and van der Goot propose that several mechanisms act at different length scales, from alignment of protein aggregates, through elongation of a dispersed phase, to separation of water at the millimeter scale~\cite{vandersman2023}. Its structure emerges from flow and cooling, which makes it continuous and scalable but gives limited control over the arrangement of phases. Fat, for example, is difficult to place in defined regions.

\section{Shear cells}

A shear cell holds a protein mixture between two surfaces, one of which rotates, while the material is heated. The shearing flow aligns the proteins, and on cooling they set in a fibrous arrangement. Shear cells, developed in academic research, allow the relationship between shear, temperature, and fiber formation to be studied under well-controlled conditions. Krintiras and colleagues produced fibrous soy-based materials in a Couette cell, a device with concentric rotating cylinders, using simple shear and heat~\cite{krintiras2015}. Shear cells produce fibrous material in batches with a single dominant direction.

\section{Spinning}

Protein fibers can be formed by forcing a protein solution through fine openings into a setting bath or through air, as synthetic textile fibers are spun. Wet spinning and related methods can produce fine, well-aligned fibers. The difficulty is assembling them into a food: fine fibers must be bundled and bound, and the throughput needed to make bulk food from individual fibers is high.

\section{Printing}

Three-dimensional printing of food deposits a paste or gel along programmed paths, layer by layer. It offers direct control over placement: fat can be put where it is wanted, and different materials can be deposited in defined patterns. Its limits are resolution and speed. A printed line is a bead of paste, not a fiber, and the fineness of the structure is set by nozzle diameter. Finer nozzles give finer structure but slower deposition, so printing whole portions with fine structure is slow.

\section{Cultivated tissue}

Cultivated meat grows animal cells in culture, sometimes on scaffolds that give them a structure to organize around. It is the only route in which structure arises from biological growth in the same way as in an animal. It is also the most technically demanding and, at present, the most expensive.

\section{Where the strand method sits}

\begin{table}[h]
\centering
\begin{tabularx}{\linewidth}{@{}lXX@{}}
\toprule
Method & What it controls & What it leaves uncontrolled \\
\midrule
High-moisture extrusion & Alignment from flow; continuous output & Placement of distinct phases \\
Shear cell & Alignment under defined shear & Scale; multiphase arrangement \\
Spinning & Fine fiber diameter and alignment & Assembly into bulk food \\
Printing & Placement of materials & Fiber fineness at useful speed \\
Cultivated tissue & Biological structure & Cost; scale \\
Core--shell strands & Phase fraction, strand size, orientation, bonding (proposed) & To be determined experimentally \\
\bottomrule
\end{tabularx}
\caption{Existing structuring methods compared with the proposed strand method.}
\end{table}

\section{Assembly and bulk structuring}

A review of structuring processes by Dekkers, Boom, and van der Goot distinguishes two broad strategies~\cite{dekkers2018}. In one, structure is built by assembling smaller structural elements, such as fibers or particles, into a larger whole. In the other, a bulk material is structured as a whole, for example by shearing or extruding it so that fibrous structure forms throughout. Spinning and printing belong mainly to the first strategy; extrusion and shear cells to the second.

The strand method belongs to the assembly strategy. It forms discrete composite elements, the strands, and then assembles them by coiling and consolidation. Positioning it this way makes clear what it inherits from earlier assembly methods, including the problem of binding many elements into a cohesive whole at useful throughput, and what it proposes to add: control of internal phase arrangement within each element, set at the preform scale and preserved by drawing.

The strand method combines features of several of these. Like spinning, it makes discrete strands. Like printing, it places distinct phases in defined positions, but it does so in the preform, at a coarse and easily controlled scale, and then reduces the scale by drawing rather than by using fine nozzles. Like extrusion, it can in principle be run continuously. Whether it achieves these advantages in practice is exactly what the experiments are for.

%----------------------------------------------------------------------
\chapter{Drawn, Pulled, and Coiled Forms}\label{ch:drawn}

Status: summarizes established crafts and industrial processes.

The individual operations of the strand method, making a composite preform, drawing it, and assembling many strands, are old. This chapter collects precedents from confectionery, noodle making, cheese making, glass working, optical fiber production, and composite manufacturing.

\section{Pulled and filled confectionery}

Hard candy is worked while warm and plastic, pulled and folded to aerate it and give it a satiny texture, and then drawn into ropes. Filled candies are made by wrapping the pulled sugar around a soft center and drawing the composite rope down to size before cutting it. The Chicken Bones candy, a cinnamon-flavored pulled sugar around a chocolate center, is a regional example. The center remains a continuous core through the whole length of the drawn rope, and its proportion to the sugar shell is set when the rope is first assembled.

Dragon's beard candy, a traditional confection, goes further. A block of hardened sugar syrup is stretched into a loop, dusted with starch to keep strands from sticking, doubled, and stretched again, repeatedly. Each doubling multiplies the number of strands by two, so a dozen or so doublings produce thousands of fine threads.

\section{Hand-pulled noodles}

Hand-pulled noodles use the same repeated doubling with wheat dough. A thick rope is stretched, folded, and stretched again, and after each cycle the number of strands doubles and their diameter falls. Gluten allows the dough to stretch without breaking and helps the strands keep their shape. The technique shows that fine, uniform food strands can be made by drawing from a coarse starting form without any fine nozzle.

\section{Stretched curd cheese}

Mozzarella and related cheeses are made by heating curd and stretching and kneading it in hot water. The stretching aligns the casein proteins into fibers, which is why string cheese peels into strands. This is a direct precedent for Chapter~\ref{ch:draw}'s observation that drawing can orient a protein-rich material and give it internal grain.

\section{Glass cane and fiber preforms}

Glassworkers make patterned cane by assembling colored glass into a thick rod and drawing it out while hot. The pattern in the cross-section shrinks uniformly while its proportions are preserved, so that intricate designs can be made at a scale where they are easy to build and then reduced to a fine rod. Millefiori beads are sliced from such canes.

Optical fiber is made the same way at industrial scale. A glass preform with a core of one composition and a cladding of another is heated at the top of a tower and drawn into a fiber less than a millimeter across, often kilometers long. The ratio of core to cladding diameter is set in the preform and preserved through the draw with high precision. The drawing process, including the heat transfer and flow in the neck-down region that governs it, is treated by Jaluria~\cite{jaluria2018}.\footnote{Optical fiber drawing preserves cross-sectional proportions far more precisely than any food process is likely to, because glass viscosity can be controlled closely with temperature. The precedent shows the principle, not the tolerance.}

\section{Rope and composite winding}

Rope is made by twisting fibers into yarns and yarns into strands, a hierarchy of assemblies whose properties depend on the twist and arrangement at each level. Filament-wound composites are made by laying resin-coated fibers onto a rotating mandrel at controlled angles, building pipes, tanks, and pressure vessels whose strength in each direction is set by the winding pattern~\cite{peters2011}. Chapter~\ref{ch:coil} borrows its geometry directly from filament winding.

\section{Shared principles}

These precedents share three principles.

\begin{enumerate}
\item \textbf{Proportions are set at the coarse scale.} Core-to-shell ratios and cross-sectional patterns are built where they are easy to see and control.
\item \textbf{Drawing reduces scale while preserving proportion,} provided the materials are worked in a suitable temperature and stiffness window.
\item \textbf{Assembly geometry sets directional properties.} Whether strands run parallel, cross, or twist determines how the assembly behaves in each direction.
\end{enumerate}

The strand method applies all three to food.

%======================================================================
\part{The Shared Platform}
%======================================================================

\chapter{Ingredient Streams and Intake}\label{ch:intake}

Status: proposed.

\section{Streams rather than recipes}

A conventional food plant is organized around products: a line for bread, a line for noodles, a line for sauces. The platform proposed here is organized around ingredient streams. A stream is a supply of one kind of material, such as a grain, a legume, a root, a vegetable, or a fruit, that enters the platform, is characterized, and is then available to any product route that can use it.

The distinction matters because it changes what happens when a supply changes. In a product-organized plant, a shortage of one ingredient stops a line. In a stream-organized platform, a route specifies required material properties, and candidate streams can be substituted where they satisfy that specification. Nominal protein, starch, or fat content is not enough: protein type, flow and setting behavior, flavor, and allergen content also matter. Locally available crops can then be substituted for one another where they meet the route's relevant specifications.

\section{Candidate inputs}

Typical inputs are cereals (wheat, oats, barley, rice, maize), legumes (peas, lentils, beans, soy), roots and tubers (potato and others), oilseeds and nuts, vegetables, and fruit. Less conventional local inputs belong in the same framework. Cattail rhizomes, for example, contain substantial starch and have been used as food in several regions. They illustrate the principle that the platform is defined by what it does with inputs, not by a fixed list of them.\footnote{Wild-harvested inputs raise their own questions of identification, contamination, and harvest sustainability, which are treated under food safety in Part VII.}

\section{Variability}

Agricultural materials vary with variety, growing conditions, harvest time, and storage. Protein content of the same grain can differ appreciably between lots; fruit sugar and acid change with ripeness; roots lose moisture in storage. A platform that treats inputs as interchangeable must measure them, because substitution only works when the relevant properties are known.

\section{Intake measurements}

Each incoming lot receives a record and a minimum set of measurements:

\begin{itemize}
\item moisture content;
\item approximate protein, fat, starch, and fiber content, measured directly or taken from supplier analysis and checked periodically;
\item particle size, for milled inputs;
\item pH, for fruit, vegetables, and anything to be fermented;
\item visual and sensory inspection for spoilage, foreign material, and off-odors;
\item for wild or uncertified inputs, identification and any contaminant screening appropriate to the source.
\end{itemize}

\section{Centrifugal profiles}

A simple characterization that requires little equipment is the centrifugal profile. A standard quantity of material is blended with a standard quantity of water, and the slurry is spun in graduated tubes at a fixed speed for a fixed time. It separates into layers, typically an oil-rich top layer, an aqueous middle layer, and a sediment of starch, protein, and fiber. The volume of each layer is recorded.

The profile is a quick fingerprint of a lot. It can show that one batch of legumes has more oil than the last, or that a fruit puree has more suspended solids. It can also be run on a reference food and a prototype side by side, as the concept drawings show.

The profile is a characterization, not a test of equivalence. Two materials with similar layer volumes can differ in protein type, starch structure, aroma, and every property that depends on structure. Blending destroys structure before the spin begins. A matching profile shows that two materials are similar in gross composition and nothing more. Chapter~\ref{ch:equivalence} returns to this.

\section{Lot records}

Every lot is given an identifier that follows its material through every downstream operation. When a product is released, its record names the lots it contains. This is the basis of the batch records described in Part VII and of any recall.

%----------------------------------------------------------------------
\chapter{Fractionation Operators}\label{ch:fractionation}

Status: the operations are established; their use in a shared platform is proposed.

An operator, in the vocabulary of this book, is a unit operation that takes one or more streams and produces one or more streams with known changes in their properties. Fractionation operators separate a material into parts.

\section{Dry milling and sieving}

Grinding reduces particle size; sieving and air classification then separate particles by size and density. Because protein and starch in some seeds occur in particles of different size, air classification of finely milled legume flour can produce protein-enriched and starch-enriched fractions without water. Dry fractionation is simpler and uses less energy than wet processing but gives less complete separation~\cite{schutyser2015}.

\section{Soaking and hydration}

Soaking softens seeds and roots, begins to release soluble components, and prepares material for wet grinding. It also reduces some undesirable compounds in legumes and can be the first step of germination or fermentation.

\section{Wet grinding and separation}

Soaked material is ground in water to a slurry. Coarse solids, mainly fiber and hull, are removed by screening or pressing through cloth. The remaining liquid can be separated further by settling or centrifugation into a fat-rich layer, an aqueous layer with soluble protein and sugars, and a sediment rich in starch. This is the domestic route to soy milk, oat milk, and nut milks, and the industrial route to starch and protein fractions.

\section{Protein precipitation}

Many plant proteins are least soluble at a particular acidity. Adjusting the pH of an aqueous extract to that point, or adding salts, causes the protein to come out of solution as a curd, which can be separated from the liquid. Tofu is made by coagulating soy milk with salts or acid. The same principle concentrates protein from other legumes.

\section{Filtration}

Filters and membranes separate by particle or molecule size, from coarse cloth that retains pulp to fine membranes that separate proteins from sugars and salts. Membrane filtration is standard in dairy processing and increasingly used for plant materials.

\section{Pressing and expression}

Mechanical pressing extracts liquid from fruit, vegetables, and oilseeds. Juice pressing leaves a fiber-rich pomace; oil pressing leaves a protein-rich cake. Both residues are streams in their own right.

\section{Fruit milking}

The phrase \emph{fruit milking} appears in the earliest statement of this program as a compact label without a definition. A working definition is adopted here.

\begin{definition}
\emph{Fruit milking} is the production of a stable, pourable suspension or emulsion from fruit, vegetable, nut, or seed tissue by disrupting the tissue, extracting its soluble and finely dispersible components into water, and removing coarse solids, so that the result can be used where milk would be used.
\end{definition}

Coconut milk, almond milk, and oat milk fall under this definition. So would liquids made from other fruits and vegetables if they can be kept from separating, which is largely a question of particle size, stabilizers, and fat content. The definition makes fruit milking a combination of operators already listed, soaking, wet grinding, screening, and emulsification, rather than a new operation. Its interest lies in widening the range of inputs from which milk-like liquids are made.

\section{Every fraction goes somewhere}

Each fractionation operation produces a stream that the target product does not use: okara from soy milk, pomace from juice, bran from flour, whey from curd. The platform must have a route for each of them, as a food ingredient, feed, fermentation substrate, or compost. Chapter~\ref{ch:closure} treats this systematically.\footnote{Industrial wet milling became economical partly by finding uses for every fraction. A small platform faces the same arithmetic with fewer outlets.}

%----------------------------------------------------------------------
\chapter{Transformation Operators}\label{ch:transformation}

Status: the operations are established; their use in a shared platform is proposed.

Transformation operators change the chemistry or physical state of a stream without separating it.

\section{Fermentation}

Fermentation uses microorganisms to transform a substrate. The main kinds relevant to the platform are:

\begin{description}[leftmargin=1.5em]
\item[Lactic acid fermentation,] by lactic acid bacteria, which acidifies and thickens, as in yogurt, sauerkraut, and sourdough;
\item[Yeast fermentation,] which produces carbon dioxide and alcohol, as in bread and brewing;
\item[Mold fermentation,] in which a fungus grows through a solid substrate, binding it and changing its flavor and texture, as in tempeh, or producing enzymes that break down starch and protein, as in koji;
\item[Bacterial solid-state fermentation,] as in natto, which produces a sticky, strongly flavored product.
\end{description}

Fermentation is a batch process with a defined history. Its record includes the substrate composition, the identity and source of the culture, the inoculation rate, and the temperature and pH over time. A fermentation is released only when its endpoint measurements, typically pH and a sensory check, fall within specification. These records are what make fermentation reproducible across batches and inputs.

A fermentation record should keep three kinds of value apart. A \emph{setpoint} is what the operator asked for, such as an incubator set to a given temperature. A \emph{measured value} is what actually occurred, such as the temperature logged in the vessel over the run, which may lag, overshoot, or drift from the setpoint. A \emph{calculated value} is derived afterward from measurements, such as the rate of acidification or a yield. A description such as ``fermented at $37^\circ$C'' collapses all three into one number and loses exactly the information needed to explain why two batches with the same setpoint came out differently. The PREFER ontology for precision fermentation makes this same division, distinguishing process control variables, process measured variables, and process calculated variables, and offers a shared vocabulary for recording them~\cite{prefer2026}.

Two kinds of fermentation should also be kept distinct. The traditional fermentations listed above use naturally occurring or conventionally selected cultures to transform a food substrate. Precision fermentation uses microorganisms, often genetically engineered production strains, to manufacture specific ingredients such as individual proteins or fats, which are then separated from the organism and used elsewhere. Precision fermentation could supply ingredient streams to the platform, and its record-keeping practices are directly useful, but it is a different activity with different regulatory standing, and products made with it should be labeled and recorded as such.

Fermentation also bears on the strand method. A fermented paste has different flow and setting behavior from an unfermented one, and a mold such as the one used in tempeh binds solid particles into a cohesive block. Whether mold growth could serve as a fusion mechanism between strands is an open question worth one trial.

\section{Enzymatic treatment}

Enzymes carry out specific transformations without living organisms. Amylases break starch into sugars, which changes sweetness and viscosity and supplies sugars for browning. Proteases break proteins into smaller pieces, which changes texture and releases savory flavor. Transglutaminase links proteins together, as discussed in Chapters~\ref{ch:shell} and~\ref{ch:consolidation}. Each is controlled by temperature, time, pH, and dose, and stopped by heating.

\section{Hydration}

Dry fractions such as flours, protein concentrates, and textured vegetable protein must be rehydrated before use. The liquid used, water, broth, or a seasoned liquid, is itself a flavor vehicle. The fried protein base described in Part I depends on hydrating textured vegetable protein in broth before frying it.

\section{Heating}

Heat gelatinizes starch, which swells and thickens; denatures proteins, which then set into gels; kills microorganisms; and drives off water. Most forming and setting operations depend on heating at a controlled rate to a controlled temperature, and the heating history of each batch is part of its record.

\section{Browning}

Frying, baking, searing, and roasting produce browning reactions at surfaces, as described in Chapter~\ref{ch:meaty}. Browning is listed separately from heating because it depends on surface conditions, low moisture and high temperature, rather than on the temperature of the bulk.\footnote{A food can be fully cooked without browning, as in boiling or steaming, or browned on the outside while still raw inside. The two operations have different controls and different records.}

%----------------------------------------------------------------------
\chapter{Forming Operators}\label{ch:forming}

Status: established operations and proposed ones are marked.

Forming operators give a material its shape and internal arrangement.

\section{Rolling hotplate compression}

Status: proposed as a platform operation. A wet paste or dough is passed between heated rollers or pressed between heated plates, which flatten it into a sheet and set it at the same time. The principle is shared by drum drying, in which a slurry dries on a heated rotating drum; by industrial flatbread and wafer ovens; and by the domestic tortilla press and waffle iron. The operator produces sheets whose thickness is set by the gap and whose texture depends on moisture, temperature, and contact time. Sheets can be eaten as flatbreads, dried and broken into crackers, or used as the shell sheet in the wrapping method of Chapter~\ref{ch:preform}.

\section{Extrusion}

Extrusion forces a paste or dough through a shaped opening. It is the forming method for pasta, many snacks, textured proteins, and high-moisture fibrous products (Chapter~\ref{ch:structuring}). On the platform, a single extruder with interchangeable dies can serve several routes.

\section{Molding}

Pastes, gels, and batters are set in molds by heating, cooling, or chemical setting. Molding gives shape but little internal structure.

\section{Printing}

Status: established technique, optional on the platform. Printing deposits material along programmed paths (Chapter~\ref{ch:structuring}). Its role on the platform is limited to products where exact placement of a small quantity of material matters more than throughput, such as decorative or layered items.

\section{Emulsification}

Mixing oil into an aqueous phase with an emulsifier, and reducing droplet size by high-shear mixing or homogenization, forms an emulsion. It is listed as a forming operator because it creates structure at the droplet scale. Its main control variable is $\psi$ (Chapter~\ref{ch:emulsion}), together with droplet size and the choice of emulsifier.

\section{Strand formation}

Status: proposed. Preform construction, shell setting, drawing, coiling, and consolidation, developed in Part IV.

%----------------------------------------------------------------------
\chapter{The Operator Graph}\label{ch:graph}

Status: proposed.

\section{Products as routes}

Once operators are defined, each product can be written as a route: an ordered sequence of operators applied to one or more streams. Different products that use the same operator share the same equipment. The collection of all routes forms a graph in which operators are nodes and routes are paths.

The graph answers a practical question. A platform with limited capital cannot build everything at once. The operators used by the most routes are the ones to build first.

\section{Example routes}

The following routes use the operator abbreviations: M, milling; S, soaking and hydration; W, wet separation; Fi, filtration; Pr, pressing; Fe, fermentation; H, heating; B, browning; E, emulsification; G, gel setting; X, extrusion; R, rolling hotplate; Mo, molding; C, coaxial preform; D, drawing; K, coiling; Q, consolidation.

\begin{table}[h]
\centering
\small
\begin{tabularx}{\linewidth}{@{}lX@{}}
\toprule
Product & Route \\
\midrule
Fried protein base (practiced) & S $\to$ B \\
Flatbread & M $\to$ S $\to$ R $\to$ B \\
Noodles & M $\to$ S $\to$ X $\to$ D $\to$ H \\
Plant milk & S $\to$ M $\to$ W $\to$ E $\to$ H \\
Cultured milk-like product & S $\to$ M $\to$ W $\to$ E $\to$ H $\to$ Fe \\
Fermented legume cake & S $\to$ H $\to$ Fe \\
High-$\psi$ spread & Pr $\to$ W $\to$ E \\
Reconstructed fruit & Pr $\to$ W $\to$ Fi $\to$ C $\to$ G $\to$ Mo \\
Strand block & M $\to$ S $\to$ W $\to$ E $\to$ C $\to$ D $\to$ G $\to$ K $\to$ Q $\to$ B \\
\bottomrule
\end{tabularx}
\caption{Illustrative product routes. Wet grinding appears as M following S.}
\end{table}

\section{Reuse counts}

Counting how many of these nine routes use each operator gives:

\begin{table}[h]
\centering
\begin{tabular}{@{}lc@{\qquad}lc@{\qquad}lc@{}}
\toprule
Operator & Routes & Operator & Routes & Operator & Routes \\
\midrule
S & 7 & B & 3 & X & 1 \\
M & 5 & Fe & 2 & R & 1 \\
W & 5 & Pr & 2 & Fi & 1 \\
E & 4 & G & 2 & Mo & 1 \\
H & 4 & C & 2 & K & 1 \\
 &  & D & 2 & Q & 1 \\
\bottomrule
\end{tabular}
\caption{Operator reuse across the example routes.}
\end{table}

Soaking, milling, wet separation, emulsification, and heating are used by most routes. They form the core of the platform and should be built first. Browning, fermentation, pressing, gel setting, coaxial preform construction, and drawing are shared by two or three routes each.

\section{What the count says about the strand method}

Coiling and consolidation, the operations specific to the strand method, are used by only one route. By the reuse criterion alone, they would be built last.

This is a real limitation and should be stated as such. The case for the strand method does not rest on reuse. It rests on the proposed ability to control a combination of structural variables, phase placement, strand geometry, orientation, and bonding, that the other routes do not control independently in the same way. The graph shows, however, that most of what the strand method needs, including milling, wet separation, emulsification, preform construction, drawing, and setting, is shared with other products. The specific additions are a coiling station and a press.\footnote{A press may itself find other uses, for example in pressing curds, cakes, and laminated doughs, which would raise its reuse count in a fuller graph.}

\section{Lineage}

The graph describes routes in general. Each actual batch traces one path through it with specific materials, and that path is its lineage. A strand block's lineage might run from two input lots, through wet separation into a protein fraction, through a particular culture and fermentation run that produced a cultured paste used as the core, through preform construction, drawing, setting, coiling, and pressing, to cooking and the release tests. Each link names the lot or run it came from and points to that step's record.

Lineage is what makes a platform debuggable. When a block fails its cross-section test, the lineage shows which input lots, which culture, and which press run it shares with blocks that passed, and which it does not. Chapter~\ref{ch:batch} describes the records that each link refers to.

\section{Using the graph}

The graph is a planning tool, not a fixed design. As routes are added or tested and found wanting, the counts change. Recording the graph alongside the batch records of Part VII keeps equipment decisions tied to what the platform actually makes.

%======================================================================
\part{The Core--Shell Strand Process}
%======================================================================

\chapter{Muscle as a Reference Structure}\label{ch:muscle}

Status: the anatomy described here is established; the comparison to composite strands is proposed.

\section{Fibers, bundles, sheaths}

Skeletal muscle is organized as a hierarchy of sheathed strands. The basic unit is the muscle fiber, a single elongated cell with many nuclei, typically tens of micrometers across and often centimeters long. Each fiber is wrapped in a thin layer of connective tissue, the endomysium. Fibers are grouped into bundles, or fascicles, each wrapped in a thicker sheath, the perimysium. The whole muscle is enclosed by a third layer, the epimysium. In cross-section the result is a set of filled compartments inside a continuous network of connective tissue, repeated at two or three scales. The structure of intramuscular connective tissue and its role in meat texture are reviewed by Purslow~\cite{purslow2005}.

Several properties of meat as food follow from this arrangement. Grain is the direction in which fibers run. Meat shreds along the planes between fascicles, where the perimysium is the weakest continuous surface. Intramuscular fat, the marbling visible in a cut of beef, is deposited largely in the perimysial spaces between bundles. Tenderness depends on the state of the fiber contents and on the amount and condition of connective tissue, which is why cuts from heavily used muscles need long moist cooking before they soften.

\section{Features a strand assembly could reproduce}

A composite strand with a soft core and a set outer layer has the same topology as a sheathed fiber: a filled compartment inside a continuous skin. If many such strands are coiled together and pressed, and if their skins fuse into a connected network (Hypothesis~\ref{h:fusion}), the cross-section of the resulting block would resemble a cross-section of muscle in three respects.

\begin{enumerate}
\item \textbf{Compartmentalization.} Core material sits in separate channels rather than in one continuous mass.
\item \textbf{Grain.} Strands run in a controlled direction, so mechanical properties differ along and across them.
\item \textbf{Weak planes.} Where fusion between strands is incomplete, the block separates along strand boundaries, giving shred.
\end{enumerate}

The hierarchy can be extended by assembling in two stages. Strands are first gathered into small bundles and given a second, thicker outer layer; the bundles are then coiled and pressed. The inner shells play the role of endomysium and the bundle shells the role of perimysium. Fat or a high-$\psi$ emulsion placed in the bundle shells, rather than in the strand cores, would place marbling where it occurs in meat.

\section{Features it probably cannot}

The comparison has limits that should be stated before any experiment is run.

\paragraph{Scale.} Food strands made by the methods in this part are likely to be a fraction of a millimeter to a few millimeters across. That is far coarser than a muscle fiber and closer to the scale of a fascicle. A strand assembly therefore resembles muscle at the bundle level, not the cellular level. This may matter less than it appears, since the shred and grain perceived in eating are largely bundle-scale phenomena.

\paragraph{Cooking contraction.} When meat is heated, its proteins denature in stages, fibers shrink across and then along their length, water is expelled, and connective tissue first tightens and, with prolonged moist heat, softens~\cite{tornberg2005}.\footnote{The temperatures at which these stages occur vary with species, muscle, and heating rate, and are treated in the food science literature as ranges rather than fixed points. The monograph does not rely on specific values.} A strand assembly will change during cooking, but not by these mechanisms, and its cooking behavior has to be measured rather than assumed.

\paragraph{Post-mortem changes.} Rigor, aging, and enzymatic tenderizing in meat have no counterpart in the proposed process.

\paragraph{Mechanical anisotropy at small strain.} Muscle fibers are themselves highly ordered protein structures. A core of fat, gel, or paste has no internal order unless drawing introduces some (Chapter~\ref{ch:draw}).

\section{What the comparison is for}

Muscle is used here as a reference structure, not a target to be matched. It provides a biological precedent for an architecture of sheathed strands in a fused network associated with grain, shred, and juiciness; it does not show that an engineered food made from the proposed materials will inherit those properties, and it supplies a vocabulary for describing what a strand block does. The question the experiments answer is not whether the block is meat but which of the three features listed above it reproduces, at what parameter values, and whether those features survive cooking.

%----------------------------------------------------------------------
\chapter{Preform Construction}\label{ch:preform}

Status: proposed.

\section{The preform}

\begin{definition}
A \emph{preform} is a composite rope, much thicker than the final strand, in which the core and outer layer are already arranged coaxially in their intended proportions.
\end{definition}

The preform is where $\phi$ is set. For a round preform of outer radius $R_0$ and core radius $r_0$,
\[
\phi = \left(\frac{r_0}{R_0}\right)^2 .
\]
A core radius of 70\% of the outer radius therefore gives $\phi \approx 0.49$, and a core radius of 90\% gives $\phi \approx 0.81$. Because $\phi$ depends on the square of the radius ratio, thin-looking shells can still hold most of the cross-section.

Building thick and drawing thin is the method of pulled candy, glass cane, and optical fiber. It avoids the need for fine coaxial nozzles, which clog and are hard to clean, and it lets the core and shell be inspected at a scale where errors are visible.

\section{Coaxial extrusion}

A coaxial nozzle feeds the core material through a central channel and the shell material through a surrounding annulus. The two streams meet at the nozzle exit. This is the natural route for continuous production and is already used for filled snacks and for co-extruded sausage casings.

Its main difficulty is rheological. When two materials of very different viscosity are extruded together, the interface does not stay where the nozzle geometry puts it; the less viscous material tends to migrate and redistribute. Core and shell formulations therefore need broadly compatible flow behavior at the extrusion temperature, or the nozzle must be designed with the mismatch in mind.

\section{Sheet wrapping}

A shell sheet is rolled or pressed flat, a core log is laid on it, and the sheet is wrapped around the core and sealed along a seam. This is how many filled confections are made; the Chicken Bones candy, for example, wraps a pulled sugar layer around a chocolate center before the rope is drawn and cut. Sheet wrapping needs no special equipment and suits bench work. Its weakness is the seam, which is a line along which the shell may split during drawing or cooking. Overlapping the seam and pressing it before drawing reduces this risk.

\section{Coating a self-supporting core}

If the core can hold its shape on its own, the shell can be applied by dipping, spraying, or enrobing. A soft core can be made temporarily self-supporting by freezing it, which is the principle behind fried ice cream: a frozen center is coated and the coating set before the center has time to melt. Frozen fat or emulsion rods can be coated in the same way. The shell must then be set and the preform drawn at temperatures the core can tolerate, which couples the choice of core material to the whole downstream sequence.

\section{Bilateral feeds}

The concept drawings show two feeds entering from opposite sides. In preform terms this can mean either of two things, and they should be kept distinct.

\begin{itemize}
\item \textbf{Two preform types.} Two separate preforms with different formulations are made and later coiled together, so that the block contains alternating strand types.
\item \textbf{Bicomponent preforms.} A single preform contains two core materials, or two shell materials, arranged side by side. This produces strands that are asymmetric in cross-section.
\end{itemize}

The first is simpler and is the one assumed in the rest of this part.

\section{Requirements on the preform}

A usable preform must hold its shape while waiting to be drawn, keep its core centered, have no seams or voids that will open under tension, and have a uniform $\phi$ along its length. Each of these is measurable by sectioning sample preforms at intervals and imaging the cut faces.

%----------------------------------------------------------------------
\chapter{Shell Setting}\label{ch:shell}

Status: proposed. This is the central engineering problem of the process.

\section{Requirements on the shell}

The outer layer of a strand has to meet three requirements that pull against one another.

\begin{description}[leftmargin=1.5em]
\item[Containment.] It must hold the core in place through drawing, coiling, and pressing, without leaking.
\item[Flexibility.] It must bend around the mandrel without cracking.
\item[Bondability.] It must still be able to fuse to neighboring shells when pressed.
\end{description}

A fully set shell should contain well but may crack and is expected to bond poorly. A lightly set shell may retain greater bondability but may fail to contain. The design problem is to find a setting route, or a sequence of routes, that passes through a state satisfying all three at the moment each is needed.

\section{Bending strain}

Flexibility can be made quantitative. When a strand of radius $R$ is bent so that its centerline has radius of curvature $\rho$, the outer surface is stretched by a strain of approximately
\[
\varepsilon \approx \frac{R}{\rho}.
\]
Coiling onto a mandrel of radius $\rho$ therefore imposes a surface strain proportional to strand thickness. Two consequences follow. Thinner strands tolerate tighter coiling, so drawing before coiling reduces the demand on the shell. And a shell can be tested for coiling in advance by bending single strands around rods of known radius and recording the smallest radius at which the shell survives.

\section{Thermal setting}

Heating the outer surface of a strand while the core stays cool sets the shell by protein gelation or starch gelatinization in a thin skin. This is the heated passage in the concept drawings. Skin thickness is controlled by the wall temperature and the residence time in the heated section. Thermal setting is continuous, needs no added chemistry, and uses the same reactions that later give cooked texture.

Its risk is that heat-set protein skins may be too stiff for coiling and may have spent the reactivity they need for fusion. A partial set, enough to contain the core but short of full gelation, is the target.

\section{Ionic setting}

A shell containing sodium alginate sets almost instantly on contact with a calcium solution, forming a gel from the surface inward.\footnote{The mechanism is the cross-linking of alginate chains by divalent calcium ions. It is the basis of the spherification technique used for fruit caviar and of calcium-set alginate sausage casings. The chemistry and gelation of alginate are reviewed by Lee and Mooney~\cite{lee2012}.} The skin thickens with contact time, so bath immersion sets the shell thickness directly. Calcium alginate gels are flexible and do not melt on heating, which helps a strand survive cooking.

The same property is a drawback for fusion: set alginate surfaces do not readily bond to each other. Ionic setting may therefore work best as a temporary outer skin over a protein-rich layer that fuses under heat and pressure, or with a binder applied at the press.

\section{Frying}

Brief frying forms a skin by rapid surface dehydration and browning. It is expected to give strong containment and adds flavor, but fried skins tend to become brittle as they cool. They are pliable while hot, so frying could be followed immediately by coiling at temperature. Frying is also the setting route most likely to spend the shell's bondability.

\section{Drying}

Surface drying, as in pasta manufacture, forms a firm skin without heat damage. Its distinctive expected property is reversibility: a dried skin might be handled and coiled, then remoistened at the press to recover its ability to bond. Drying may therefore allow the containment and fusion stages to be separated, which the other routes do less directly. Its cost is time, and uneven drying can crack the skin.

\section{Enzymatic setting}

Transglutaminase forms covalent cross-links between protein chains and is already used to bind meat and fish pieces~\cite{kuraishi2001}. Included in the shell formulation, it sets the skin slowly over minutes to hours. The slowness is useful if the strand is coiled and pressed while cross-linking is still under way, so that fusion and setting happen together. It requires close control of timing and temperature.

\section{Comparison}

\begin{table}[h]
\centering
\begin{tabularx}{\linewidth}{@{}lXXX@{}}
\toprule
Route & Containment & Flexibility & Bondability \\
\midrule
Thermal & Good if skin forms & Uncertain; may be stiff & Reduced by full set \\
Ionic & Good & Good & Poor without binder \\
Frying & Strongest & Good hot, poor cold & Likely poor \\
Drying & Good & Moderate & Recoverable by remoistening \\
Enzymatic & Slow to develop & Good early & Good if pressed in time \\
\bottomrule
\end{tabularx}
\caption{Expected behavior of setting routes. Every entry is a prediction to be tested in Experiment E3.}
\end{table}

\section{Sequence}

Setting does not have to be a single step. A strand could be given a light ionic skin for handling, drawn, coiled, and then heat-set and fused in the press. Or it could be skin-dried, coiled, remoistened, and pressed. The choice of sequence determines where drawing can occur, because a fully set shell will not draw. Chapter~\ref{ch:draw} returns to this.

\section{Tests}

Each setting route is judged by three measurements on single strands: leakage of core material under light pressure; the minimum bend radius survived without cracking; and, after two strands are pressed together under standard conditions, the force needed to peel them apart. These are the three requirements turned into numbers.

%----------------------------------------------------------------------
\chapter{Draw-Down}\label{ch:draw}

Status: proposed.

\section{Volume conservation}

When a preform of diameter $D_0$ and length $L_0$ is drawn to diameter $D = D_0/\lambda$, conservation of volume requires its length to become
\[
L = \lambda^2 L_0 .
\]
If core and shell stretch together with the same axial strain, with no slip at the interface, every cross-sectional area is multiplied by the same factor $1/\lambda^2$. The ratio of core area to total area is then unchanged:
\[
\phi_{\text{strand}} = \phi_{\text{preform}}.
\]
This is the content of Hypothesis~\ref{h:draw}. A draw ratio of 10 turns a 20\,mm preform into a 2\,mm strand one hundred times longer, with the same core fraction.

\section{When the invariance fails}

The idealization breaks down in several ways, each of which is a thing to look for in cross-sections.

\paragraph{Viscosity mismatch.} If the core is much softer than the shell, the shell carries nearly all the drawing force. The core may then lag, bunch, or thin unevenly. If the core is much stiffer, it may resist drawing and tear the shell.

\paragraph{Capillary breakup.} A liquid core in a thin, soft shell can become unstable as the diameter falls: surface tension drives it to break into a chain of droplets, as a thin stream of water breaks into drops. This would turn a continuous core into a beaded one. That may be a defect or a feature; a beaded core resembles fat pockets more than fat channels.

\paragraph{Necking.} Many materials draw unevenly, forming thin necks that break. Stable drawing requires a material that stiffens as it stretches, which wheat doughs do because of gluten, or careful temperature control, as in pulled candy and glass, which are drawn in a window where they are soft but not runny.

\paragraph{Seam opening.} A wrapped preform may split along its seam.

\section{Orientation from drawing}

Drawing does more than thin the strand. Extension tends to align long protein and starch molecules along the draw direction, which is one of the ways extrusion produces fibrous texture. A drawn shell may therefore be anisotropic in its own right, adding fine-scale grain within each strand. This is a possible benefit that should be measured rather than assumed, for example by comparing the tensile strength of shell material along and across the draw direction.

\section{Where drawing fits in the sequence}

A fully set shell will not draw. There are two consistent sequences.

\begin{enumerate}
\item \textbf{Draw, then set.} The preform is drawn while both phases are soft, and the shell is set on the fine strand, for example in a heated passage or setting bath. This is simpler but asks the soft preform to survive drawing unaided.
\item \textbf{Light set, draw, full set.} The preform is given a thin, drawable skin, drawn, and then set further. This gives more control over the core during drawing but requires a skin that stretches without tearing.
\end{enumerate}

\section{Measuring \texorpdfstring{$\phi$}{phi} after drawing}

Samples are cut at intervals along the strand and the cut faces photographed under magnification. Core and shell are distinguished by color, which can be added to the core as a tracer for this purpose, and their areas measured by thresholding the image. The mean and spread of $\phi$ along the strand, compared with the preform value, test Hypothesis~\ref{h:draw} directly.

%----------------------------------------------------------------------
\chapter{Coiling Geometry}\label{ch:coil}

Status: proposed.

\section{Helix angle and pitch}

Strands are laid onto a rotating mandrel of radius $\rho$ while a guide traverses along its axis. Each strand follows a helix. Measuring the helix angle $\theta$ from the mandrel axis, the axial advance per turn, or pitch, is
\[
p = \frac{2\pi\rho}{\tan\theta}.
\]
When $\theta$ approaches $90^\circ$ the pitch approaches zero and strands lie side by side around the circumference; this is hoop coiling. A single strand of diameter $d$ laid as a closed layer requires $p = d$. A band of $k$ strands laid together requires $p = kd$, which allows lower angles at the same coverage.

\section{Crossed layers}

Reversing the traverse direction lays the next layer at $-\theta$, so successive layers cross. This is how filament-wound composite pipes are built, and it produces a laminate whose layers resist separating in the directions the strands run. In food terms, parallel layers are expected to favor shred and crossed layers to favor resistance. The angle between layers is a continuous variable, so the transition between them can be mapped.

\section{Bilateral coiling}

Two strand bands fed from opposite sides of the mandrel can be laid simultaneously at $+\theta$ and $-\theta$, producing crossed layers in a single pass. Alternatively, two strand types, such as a lean strand and a fat-rich strand, can be laid in alternating bands, intended to produce a marbled pattern when the block is cut.

\section{The central bore}

Removing the mandrel leaves a hollow center. There are four ways to deal with it.

\begin{enumerate}
\item \textbf{Edible mandrel.} Coil onto a log of cooked dough or pressed strands, which stays in the product.
\item \textbf{Flatten.} Remove the mandrel and press the hollow cylinder flat. The bore closes and the coil becomes a two-sided slab in which strands run in the plane of the slab. This is the press stage shown in the concept drawings.
\item \textbf{Slit and unroll.} Cut the coil along its length and open it into a flat mat of parallel strands at angle $\theta$ to the cut edge. Mats can then be stacked at chosen angles, as composite laminates are, giving direct control over the orientation of each layer.
\item \textbf{Slice lengthwise.} Keep the cylinder and cut it into long pieces that exclude the bore.
\end{enumerate}

Slitting and unrolling is the most general, because it separates the making of oriented mats from the design of the final layup.

\section{Grain in the finished piece}

A coiled cylinder has three directions: axial, circumferential, and radial. With hoop coiling the strands run circumferentially, so grain is circular around the axis. How the grain meets the teeth depends on how the piece is cut. A disc cut across the axis presents strands lying in the plane of the disc, while a slab cut from a flattened coil presents them running across the slab. The product map in Part V has to specify cut orientation as well as coiling angle, because the same block cut two ways will eat differently.

\section{Tension}

Strands are laid under some tension, which helps them pack but also squeezes their cores and strains their shells. Tension adds to the bending strain of Chapter~\ref{ch:shell} and should be kept low and constant, set by a tensioning guide rather than by hand.

%----------------------------------------------------------------------
\chapter{Consolidation}\label{ch:consolidation}

Status: proposed.

\section{Packing and voids}

Round strands laid side by side do not fill space. In the densest arrangement, hexagonal packing, circles cover a fraction
\[
\eta_{\text{hex}} = \frac{\pi}{2\sqrt{3}} \approx 0.907
\]
of the cross-section. In square packing they cover $\eta_{\text{sq}} = \pi/4 \approx 0.785$. The remaining 9 to 22\% is void. Before pressing, the core fraction of the whole block is therefore not $\phi$ but $\phi\eta$.

Pressing deforms round strands toward polygons that fill the voids. If core and shell deform in proportion and neither is squeezed out, a fully densified block has core fraction $\phi$, and in hexagonal packing each shell becomes a hexagonal wall shared with its neighbors. That shared-wall network is the honeycomb cross-section described in Chapter~\ref{ch:muscle}, and it exists only if the walls fuse.

\section{Fusion mechanisms}

Adjacent shells can be joined in several ways.

\begin{description}[leftmargin=1.5em]
\item[Co-gelation.] Protein-rich surfaces pressed together and then heated gel as one, as in sausage and surimi products.
\item[Remoistening.] Dried skins are misted before pressing and bond as they rehydrate.
\item[Enzymatic cross-linking.] Transglutaminase, in the shell or applied to the surfaces, links proteins across the interface.\footnote{Transglutaminase forms bonds between glutamine and lysine residues on protein chains. Its effectiveness depends on the protein source, since some plant proteins offer fewer accessible sites than meat or dairy proteins.}
\item[Binders.] A thin protein slurry or starch paste between layers acts as glue.
\end{description}

\section{Pressure, temperature, and time}

Pressure closes voids and brings surfaces into contact. Too much ruptures shells and expels cores. Temperature drives fusion but also softens or melts cores, which makes rupture more likely under load. This suggests a two-stage press: compress cold until voids close, then heat under a lower holding pressure until the shells fuse. Dwell time at temperature sets how far fusion proceeds.

\section{The fusion index}

\begin{definition}
The \emph{fusion index} $f$ is the ratio
\[
f = \frac{F_{\text{sep}}}{F_{\text{ref}}},
\]
where $F_{\text{sep}}$ is the peak force needed to pull apart two strands pressed together over a standard bonded length, and $F_{\text{ref}}$ is the peak force needed to break a single strand of the same type, given the same press and heat treatment, in tension. Both forces are measured with the same clamps and pulling rate. The two specimens necessarily differ in geometry, a bonded pair pulled apart from its tails and a single strand pulled along its length, so each geometry is specified separately and both are reported with the result (Protocol~\ref{prot:peel}).
\end{definition}

The fusion index is an operational ratio of two forces. It is dimensionless and useful for comparing conditions within one test geometry, but it depends on that geometry and is not an intrinsic bond strength of the interface. Every reported value is accompanied by both specimen geometries, the bonded length, and the location of failure in each test: at the interface, within a shell, or through a ruptured core. When a pair breaks through a strand before the interface separates, the result is recorded as $f \geq 1$ with cohesive failure.

The working prediction is that low $f$ will favor clean separation between strands and therefore shredding, while values near or above 1 with cohesive failure will correspond to conditions in which the interface is no longer the weakest surface and the block behaves more nearly as a continuous solid with embedded cores. Intermediate values may produce intermediate textures. Experiment E6 tests whether these mechanical regimes correspond to those separation and eating behaviors; the definition of $f$ does not by itself calibrate it.

The fusion index may prove to be a more direct control of shred than coiling angle. Coiling sets the directions along which a block can separate; fusion decides whether it does. A block of hoop-coiled strands with high $f$ may eat like a dense slab, while the same geometry with low $f$ pulls apart. Mapping texture over the pair $(\theta, f)$ is therefore a central task of Experiment E6.

\section{Oil and fusion}

Oil or seasoning applied between layers during coiling, as shown in the concept drawings, lowers fusion wherever it lies. That is a constraint and also a tool: a deliberate pattern of oiled interfaces creates planned weak planes, reproducing the role of the perimysium as the surface along which meat divides.

\section{Measuring consolidation}

Pressed blocks are sectioned before and after cooking. Images of the cut faces show void closure, shell wall continuity, and core integrity. Peel tests on pressed strand pairs give $f$. Core material found outside its channels after pressing or cooking measures rupture. Together these test Hypothesis~\ref{h:fusion}.

%----------------------------------------------------------------------
\chapter{Finishing}\label{ch:finishing}

Status: seasoning practice is practiced; its application to strand blocks is proposed.

\section{Three seasoning sites}

A strand block offers three places to put flavor, and each behaves differently.

\paragraph{In the core.} Seasoning mixed into the core material is distributed through the whole product before it is formed, which would address the problem noted in the concept drawings on seasoning, where surface seasoning does not reach the center. Many aroma compounds dissolve in fat, so a fat or emulsion core could serve as a carrier for them and release them during cooking and chewing. How much is released, and when, depends on core composition, emulsion stability, and temperature, and has to be measured.

\paragraph{Between layers.} Seasoning applied during coiling places flavor along interfaces. As noted in Chapter~\ref{ch:consolidation}, anything placed between layers reduces fusion there, so interlayer seasoning and texture have to be designed together.

\paragraph{On the surface.} Surface seasoning and browning form the crust, which is the first thing smelled and tasted and often dominates the perceived identity of the food.

\section{Browning}

Browning requires a surface that is hot, relatively dry, and supplied with sugars and amino acids. These can be built into the shell formulation, so that the outer strands of a block brown readily when seared. The block interior, protected by its outer layers, stays moist. This is the same arrangement as a seared piece of meat, produced by composition rather than by anatomy.

\section{Identity at the last stage}

If Hypothesis~\ref{h:late} holds, the finishing profile $F$ decides what a block is taken to be. The same strand block could be finished toward the aroma profile of a familiar meat, of a vegetable or legume dish, or of nothing with an established name, by changing aroma, fat, salt, acid, and bitterness at this stage alone. Finishing is also where bitter notes such as coffee or cocoa, and the aroma of fried onion, can supply the roasted and savory character associated with seared meat. The block itself would then be a general texture, and product identity a matter of the last few operations. This is the strongest practical consequence of the hypothesis and the main reason to test it early.

\section{Portioning and cooking}

Portions are cut with a stated orientation relative to the grain, since cut direction changes eating behavior. During cooking, fat cores melt and give juiciness, shells may continue to set, and moisture is lost. Cooking loss, measured as the weight lost on a standard cooking procedure, and the condition of the cross-section after cooking, are recorded for every block. A block whose internal structure disappears on cooking has failed regardless of how it looked raw. Safe cooking and holding procedures are set out in Part VII.

%======================================================================
\part{Specifications}
%======================================================================

\chapter{Parameter Set}\label{ch:params}

Status: proposed.

\section{Control parameters}

\begin{table}[h]
\centering
\begin{tabularx}{\linewidth}{@{}llX@{}}
\toprule
Symbol & Parameter & Proposed role \\
\midrule
$\psi$ & Emulsion phase fraction & Consistency of emulsion cores and spreads \\
$\phi$ & Core area fraction & Proportion of core material; juiciness and fat release \\
$d$ & Strand diameter & Fineness of grain \\
$\lambda$ & Draw ratio & Links preform to strand \\
$\theta$ & Helix angle & Orientation of strands in the block \\
$n$ & Layer count & Block thickness \\
$P,T,t$ & Press conditions & Void closure and fusion \\
$f$ & Fusion index & Shred versus continuous solid \\
$C,S$ & Core and shell formulations & Phase chemistry \\
$F$ & Finishing profile & Aroma, surface, perceived identity \\
\bottomrule
\end{tabularx}
\caption{Control parameters. All roles are hypotheses until measured.}
\end{table}

\section{Set parameters and resulting parameters}

The parameters are not all independent. Some are set directly by the operator; others result from those settings. Keeping the distinction clear matters for experimental design, because only set parameters can be varied at will.

\begin{description}[leftmargin=1.5em]
\item[Set directly:] preform radii $R_0$ and $r_0$; draw ratio $\lambda$; formulations $C$ and $S$; setting route and its time and temperature; mandrel radius $\rho$ and helix angle $\theta$; layer count $n$; press conditions $P$, $T$, $t$; finishing profile $F$.
\item[Resulting:] core fraction $\phi = (r_0/R_0)^2$, subject to Hypothesis~\ref{h:draw}; strand diameter $d = 2R_0/\lambda$; packing fraction $\eta$ before pressing; fusion index $f$, which depends on $S$, the setting route, and $P$, $T$, $t$; and the measured texture and anisotropy of the block.
\end{description}

\section{Couplings}

Several relations constrain which combinations are possible.

\paragraph{Bending.} Coiling imposes a surface strain $\varepsilon \approx R/\rho = d/(2\rho)$ on each strand. If the shell survives strains up to $\varepsilon_{\max}$, measured in the bend test, then the mandrel radius must satisfy
\[
\rho \geq \frac{d}{2\,\varepsilon_{\max}} .
\]
Thicker strands or more brittle shells require larger mandrels.

\paragraph{Pitch.} For a closed single-strand layer, $2\pi\rho/\tan\theta = d$, so helix angle, mandrel radius, and strand diameter are linked. Laying bands of $k$ strands relaxes this to $kd$.

\paragraph{Pressing.} Higher press temperature raises $f$ but softens the core, lowering the pressure at which cores rupture. Temperature and pressure therefore trade against each other.

\paragraph{Setting and drawing.} A fully set shell cannot be drawn. The setting route determines where drawing can occur in the sequence (Chapter~\ref{ch:draw}).

\section{First trial matrix}

Initial ranges for these parameters are not asserted here; they are to be found. The first trials use a small number of levels chosen to span a wide range, so that effects, if present, are large enough to see:

\begin{center}
\begin{tabular}{@{}ll@{}}
\toprule
Parameter & Levels \\
\midrule
$\phi$ & 0, 0.25, 0.5, 0.75 \\
$d$ & approximately 1, 2, and 4\,mm \\
$\theta$ & near-hoop (single direction); crossed at about $\pm45^\circ$ \\
Setting route & dried skin; thermal skin; ionic skin \\
$f$ & low (oiled interfaces); high (binder or co-gelation) \\
\bottomrule
\end{tabular}
\end{center}

These are levels for an experiment, not recommended values. A full factorial across them would be large. Chapter~\ref{ch:sequence} orders the work so that each stage fixes some parameters before the next stage varies others.\footnote{Fixing the shell formulation and setting route first, on the basis of single-strand tests, reduces the later block experiments to the geometric parameters $\phi$, $d$, $\theta$ and the fusion level.}

%----------------------------------------------------------------------
\chapter{The Product Map}\label{ch:map}

Status: every region on the map is a hypothesis.

\section{Structure first}

The product map places candidate foods according to their structure rather than according to the foods they might resemble. Its axes are core fraction $\phi$ and strand diameter $d$. Two further variables, helix angle $\theta$ and fusion index $f$, strongly affect texture and are held fixed on any single map; a full description would be a set of maps, one for each combination.

\begin{figure}[h]
\centering
\begin{tikzpicture}[x=2.1cm,y=7cm]
\draw[->] (0,0) -- (4.3,0); \node[below] at (2.1,-0.07) {strand diameter $d$ (mm, logarithmic)};
\draw[->] (0,0) -- (0,0.95) node[above] {core fraction $\phi$};
\foreach \x/\l in {0/0.25,1/0.5,2/1,3/2,4/4} \draw (\x,0) -- (\x,-0.015) node[below] {\l};
\foreach \y in {0,0.25,0.5,0.75} \draw (0,\y) -- (-0.05,\y) node[left] {\y};
\fill[gray!20] (0,-0.005) rectangle (4.2,0.03);
\node[font=\small] at (1.4,0.055) {noodles, pasta, solid strands};
\draw[dashed,rounded corners] (0.1,0.12) rectangle (1.4,0.42); \node[font=\small,align=center] at (0.75,0.27) {fine shred};
\draw[dashed,rounded corners] (1.6,0.12) rectangle (2.9,0.42); \node[font=\small,align=center] at (2.25,0.27) {coarse shred};
\draw[dashed,rounded corners] (3.0,0.11) rectangle (4.2,0.32); \node[font=\small,align=center] at (3.6,0.215) {dried chewy\\strips};
\draw[dashed,rounded corners] (0.6,0.46) rectangle (2.6,0.66); \node[font=\small,align=center] at (1.6,0.56) {dense sliceable block\\(high fusion)};
\draw[dashed,rounded corners] (2.8,0.45) rectangle (4.2,0.88); \node[font=\small,align=center] at (3.5,0.665) {reconstructed\\fruit and\\filled strands};
\draw[dashed,rounded corners] (0.1,0.70) rectangle (2.6,0.88); \node[font=\small,align=center] at (1.35,0.79) {soft, juicy, fat-releasing};
\end{tikzpicture}
\caption{Hypothesized regions of the product map at a single helix angle. Dashed outlines are guesses to be replaced by measured boundaries. The region labeled ``dense sliceable block'' assumes high fusion; at low fusion the same $\phi$ and $d$ would be expected to shred.}
\end{figure}

\section{Regions}

\paragraph{$\phi = 0$.} Strands with no core are ordinary solid strands. With no coiling they are noodles; coiled and pressed with low fusion, they are a nest or cake of noodles; with high fusion, a dense block of solid material with grain.

\paragraph{Low $\phi$, fine strands.} Thin strands with small cores and relatively thick shells should give a firm, fibrous material that shreds finely if fusion is low. This is the region in which structures comparable to shredded poultry are expected, if they are found at all.

\paragraph{Low $\phi$, coarse strands.} The same composition in thicker strands should shred more coarsely.

\paragraph{Coarse strands, dried.} Thick strands with low $\phi$, dried after consolidation, should give chewy strips.

\paragraph{Intermediate $\phi$, high fusion.} Cores make up half or more of the cross-section and shells are well fused. The block should slice cleanly and resist tearing, with grain visible on the cut face.

\paragraph{High $\phi$.} Thin shells around large soft cores should give soft, juicy blocks that release fat or liquid on cooking and chewing, provided the shells contain the cores through cooking.

\paragraph{Coarse, high-$\phi$ strands with gel cores.} Large gel cores in thin skin-like shells resemble the structure of some fruits, which consist of juice-filled cells or segments in thin membranes. Citrus segments are the clearest example. This region is where reconstructed fruit would be made.

\paragraph{Confectionery.} Filled candies occupy the same map with sugar shells. They are not a target of this book but mark the region where the method is already known to work.

\section{Reading the map correctly}

The map is a set of predictions, not results. Its boundaries are drawn where a first guess puts them. Experiment E7 replaces each dashed outline with measured boundaries or removes it. A region that turns out to be unreachable, for example because shells cannot contain cores at high $\phi$ through cooking, is as useful a result as a confirmed one.

\section{Emulsion products}

A companion map for emulsion products has a single structural axis, $\psi$, with droplet size and emulsifier as held variables. Along it lie pourable milks, creams, spoonable sauces, and spreads. Emulsions also appear on the strand map as core materials, so a high-$\psi$ emulsion core links the two maps.\footnote{The emulsion map is included to show how the platform's emulsification operator serves products outside the strand method, not as a claim that $\psi$ determines emulsion texture by itself.}

%----------------------------------------------------------------------
\chapter{Reference Formulations}\label{ch:formulations}

Status: proposed starting points. None has been tested in the strand process.

The formulations below are starting points for trials, chosen because each uses widely available ingredients and a known setting behavior. Proportions are given as approximate ratios to be adjusted in the first trials.

\section{Shell formulations}

\paragraph{S1: Wheat gluten dough.} Vital wheat gluten hydrated with roughly an equal weight of water or broth, with a small amount of flour or starch to adjust stiffness. Gluten dough stretches without breaking, which suits drawing, and sets on steaming or baking. It is the shell of choice for first drawing trials. It excludes people who cannot eat gluten, so a second shell is needed for any general product.

\paragraph{S2: Protein paste with alginate.} A legume protein isolate or concentrate hydrated to a stiff paste, with sodium alginate added at a level on the order of one percent of the paste weight, set in a calcium bath of similar strength. The alginate gives an immediate flexible skin; the protein beneath it sets later on heating and is expected to carry fusion.

\paragraph{S3: Starch and protein dough for drying.} A dough of rice or wheat flour with added protein, formed into a skin and surface-dried. This shell tests the dried-skin route, in which bonding is recovered by remoistening at the press.

\section{Core formulations}

\paragraph{C1: Fat core.} A fat that is solid at room temperature and melts in the mouth, such as refined coconut oil, blended with fat-soluble seasoning. It can be frozen for handling and coated by dipping. It is the simplest way to test whether shells contain a melting core through cooking.

\paragraph{C2: High-$\psi$ emulsion core.} Oil emulsified into a plant milk or cooking liquid at a high oil fraction, as in egg-free mayonnaise. It carries fat and aroma in a spoonable form that holds its shape.

\paragraph{C3: Gel core.} Fruit puree or broth set with a gelling agent such as agar, which sets on cooling and remains solid at temperatures well above body temperature. Gel cores are used for reconstructed fruit and for testing containment of a watery phase.

\paragraph{C4: Seasoned protein paste core.} Textured vegetable protein hydrated in seasoned broth, fried briefly with onion and herbs, then minced and bound with a small amount of starch. This places the practiced fried protein base of Part I at the center of a strand, so that the core carries a known and tested flavor.

\section{Pairings for first trials}

\begin{center}
\begin{tabular}{@{}lll@{}}
\toprule
Trial & Shell & Core \\
\midrule
Drawing and orientation & S1 & C1, C4 \\
Ionic skin and containment & S2 & C1, C3 \\
Dried skin and remoistened fusion & S3 & C4 \\
\bottomrule
\end{tabular}
\end{center}

Each pairing isolates one question. The gluten shell with fat and paste cores asks whether drawing preserves $\phi$. The alginate shell with fat and gel cores asks whether an ionic skin contains soft cores. The dried shell with paste core asks whether a dried skin can be remoistened to fuse.\footnote{Formulations should be recorded by weight to the nearest gram, with lot identifiers for every ingredient, so that a successful trial can be repeated.}

%----------------------------------------------------------------------
\chapter{Bench Apparatus}\label{ch:bench}

Status: proposed. Every item is ordinary kitchen or workshop equipment or a simple modification of it.

\section{Preform construction}

\begin{itemize}
\item \textbf{Wrapping:} a rolling pin, a flat work surface lined with baking paper, and a sharp blade. The shell is rolled to a measured thickness between guide sticks of known height; the core log is rolled to a measured diameter.
\item \textbf{Coaxial piping:} two piping bags, one nested inside the other, or a purpose-made coaxial tip. Core material fills the inner bag and shell material the outer.
\item \textbf{Dipping:} a freezer, a rod mold for freezing cores, and a tall container for the shell material or setting bath.
\end{itemize}

\section{Setting}

\begin{itemize}
\item a shallow tray for ionic setting baths, with a timer;
\item a steamer or water bath held at a measured temperature, for thermal skins;
\item a rack with airflow, for surface drying.
\end{itemize}

\section{Drawing}

Drawing at the bench is done by hand, pulling the preform steadily between the hands or over a smooth rod, or between two pairs of rollers turned at different speeds. Diameter is checked with calipers at intervals along the strand.

\section{Coiling}

A rotisserie spit or a hand-cranked dowel held in two supports serves as a mandrel. A food-safe sleeve of baking paper or silicone allows the coil to be slid off. The traverse is guided by hand along a marked scale, so that pitch, and hence helix angle, can be kept approximately constant.

\section{Pressing}

A heated sandwich press or panini grill provides heated plates. For controlled pressure, a lever press or a flat plate loaded with known weights is used, with pressure computed from the load and the plate area. A probe thermometer records the block's internal temperature.

\section{Measurement}

\begin{itemize}
\item a set of dowels of graded diameter, for bend tests;
\item a spring scale or luggage scale with a clamp, for peel tests;
\item a phone camera with a macro lens and a ruler in frame, for cross-section images;
\item a kitchen scale reading to one gram, for formulations and cooking loss;
\item graduated centrifuge tubes and a small centrifuge, for profiles.
\end{itemize}

\section{Food contact and cleaning}

All surfaces in contact with food are stainless steel, food-grade silicone, food-safe plastics, or baking paper. Wooden dowels, if used, are sleeved. Equipment is washed and sanitized between trials, and allergens, especially gluten, are separated by cleaning or by keeping dedicated tools.\footnote{Bench trials are for evaluation, and samples served to tasters are prepared under the food safety practices of Chapter~\ref{ch:safety}.}

%----------------------------------------------------------------------
\chapter{Pilot Apparatus}\label{ch:pilot}

Status: proposed.

\section{Line layout}

A continuous pilot line follows the bench sequence with each step mechanized:

\begin{enumerate}
\item two metering pumps feeding core and shell materials to a coaxial die;
\item a setting stage, either a heated tube, a bath with a conveyor, or a drying tunnel;
\item draw rollers running faster than the extrusion speed;
\item a coiling station with a driven mandrel and a synchronized traverse;
\item a heated platen press with pressure and temperature control;
\item a finishing station for surface seasoning and searing.
\end{enumerate}

\section{Setting \texorpdfstring{$\phi$}{phi} continuously}

With coaxial extrusion, $\phi$ is set by the ratio of volumetric flow rates of core and shell. If the core pump delivers $Q_c$ and the shell pump $Q_s$, then, provided both leave the die at the same speed,
\[
\phi = \frac{Q_c}{Q_c + Q_s}.
\]
Changing pump speeds changes $\phi$ without changing the die, which makes $\phi$ a continuous control on the pilot line.

\section{Draw by speed ratio}

If the strand leaves the die at speed $v_1$ and is taken up by rollers at speed $v_2$, conservation of volume gives an area reduction of $v_2/v_1$, so the diameter draw ratio is
\[
\lambda = \sqrt{v_2/v_1}.
\]
A diameter reduction by a factor of four requires the take-up speed to be sixteen times the exit speed.

\section{Synchronized coiling}

For a mandrel of radius $\rho$ turning at $N$ revolutions per second and a desired helix angle $\theta$, the traverse must move at
\[
v_t = N \cdot \frac{2\pi\rho}{\tan\theta}.
\]
Synchronizing the traverse to the spindle by gearing or by electronic control keeps $\theta$ constant.

\section{Control points}

Each of the following is measured and recorded for every batch: core and shell pump rates; die and setting temperatures; setting residence time; draw speeds; coiling tension, spindle speed, and traverse speed; press pressure, temperature, and dwell; block internal temperature at the end of pressing.

\section{Throughput}

Throughput is set by strand cross-section and line speed. As an illustration, a single strand of 2\,mm diameter has a cross-sectional area of about 3.1\,mm$^2$. At a take-up speed of 0.5\,m/s it carries about 1.6\,cm$^3$ of material per second, roughly 5 to 6\,kg per hour at a density near that of water. Running many strands in parallel from a multi-orifice die multiplies this. These figures are illustrative arithmetic, not performance claims; actual speeds will be limited by setting time and by how fast strands can be drawn without breaking.\footnote{Setting time is likely to be the limiting factor. A strand moving at 0.5\,m/s through a setting stage that requires ten seconds needs a five-meter path, which argues for baths or tunnels in which the strand loops.}

%======================================================================
\part{Evaluation}
%======================================================================

\chapter{Instrumental Texture}\label{ch:texture}

Status: established methods applied to proposed products.

\section{Texture profile analysis}

Sensory texture vocabulary was standardized in the early 1960s, when Szczesniak, Brandt, and Friedman developed rating scales for mechanical parameters such as hardness, chewiness, and gumminess, each anchored by reference foods, and related them to instrumental measurements~\cite{szczesniak1963}. The instrumental methods below descend from that work, and its reference scales are the model for the descriptive panels of Chapter~\ref{ch:sensory}.

Texture profile analysis compresses a sample of standard size twice in succession with a flat plunger, imitating two bites, and records force against time. From the curve are derived hardness (peak force in the first compression), cohesiveness (ratio of work in the second compression to work in the first), springiness (how far the sample recovers between compressions), and chewiness (a product of these). The method requires a texture analyzer, a standard laboratory instrument, but its logic can be approximated at the bench with a lever, a scale, and a ruler.

For strand blocks, samples are cut as cubes with one face perpendicular to the strand direction, and the compression direction is recorded. Compressing along the strands and across them gives two different profiles, and the difference is itself a result.

\section{Shear force}

Meat science commonly measures toughness by the force needed to shear a cylindrical or rectangular sample with a blade, cutting across the fibers. The same method applied to strand blocks allows direct comparison with reference foods measured the same way.

\section{The anisotropy index}

\begin{definition}
The \emph{anisotropy index} $\mathcal{A}$ is the peak shear force for a cut parallel to the strands divided by the peak shear force for a cut across them.
\end{definition}

A material with no grain has $\mathcal{A} \approx 1$. A material that separates easily along its grain but resists cutting across it has $\mathcal{A} < 1$. The index summarizes in one number the property that grain produces, and it is the main mechanical outcome of the structural prediction stated at the start of the book: changing $\phi$ and $\theta$ should change $\mathcal{A}$.\footnote{Some texture studies define anisotropy as the inverse ratio. The convention used here should be stated wherever results are reported, so that values can be compared.} Recent work on plant-based meat analogues has measured anisotropy at several scales, combining microscopy and X-ray scattering of fiber orientation with tensile and dynamic mechanical tests along and across the fibers, and reporting anisotropy indexes as ratios of longitudinal to transverse properties~\cite{zink2024}. The shear-based index used here is a simpler counterpart suited to bench equipment; where a texture analyzer is available, directional tensile tests of that kind are the better measure.

\section{Cooking loss and water holding}

Samples are weighed before and after a standard cooking procedure. The fractional weight loss measures water and fat released. Water holding capacity is measured by centrifuging a sample under standard conditions and weighing the liquid expelled. For strand blocks with fat or emulsion cores, cooking loss also indicates whether cores have leaked.

\section{Standardization}

Every measurement specifies sample size, orientation relative to the strands, temperature at test, and the cooking procedure, if any, that preceded it. Without these, texture numbers cannot be compared between batches, let alone between laboratories.

%----------------------------------------------------------------------
\chapter{Structure and Composition}\label{ch:structure}

Status: established methods applied to proposed products.

\section{Cross-section imaging}

The most direct evidence about a strand block is a look at its cut face. Blocks are sectioned with a sharp blade perpendicular to the strands and photographed under magnification with a scale in the frame, both raw and after cooking. Color tracers in the core, such as a small amount of turmeric, beet, or cocoa, make cores easy to distinguish from shells.

From each image are measured:

\begin{itemize}
\item \textbf{core fraction}, by thresholding core color and dividing core area by total area, compared with the preform value (Hypothesis~\ref{h:draw});
\item \textbf{void fraction}, the area of gaps between strands, which indicates how completely pressing densified the block;
\item \textbf{wall continuity}, whether shell boundaries between neighboring strands are visible as lines or have merged, which indicates fusion (Hypothesis~\ref{h:fusion});
\item \textbf{core integrity}, whether cores remain in their channels or have spread into walls and voids.
\end{itemize}

Comparing raw and cooked sections from the same block shows whether structure survives cooking. A block whose raw section shows clean cores and fused walls but whose cooked section shows a uniform mass has failed the structural prediction.

\section{Composition}

Proximate analysis, which measures moisture, protein, fat, carbohydrate, and ash, is carried out on finished products by standard methods, usually through a commercial laboratory. At the bench, moisture is measured by drying to constant weight and fat content is estimated from formulation.

\section{Centrifugal profiles of products}

The centrifugal profile of Chapter~\ref{ch:intake} can be run on finished products and on reference foods. It shows gross compositional similarity or difference. It cannot show structural similarity, because blending destroys structure before the spin.

%----------------------------------------------------------------------
\chapter{Sensory Methods}\label{ch:sensory}

Status: established methods.

\section{Discrimination: the triangle test}

In a triangle test, a taster receives three coded samples, two identical and one different, and must identify the odd one. By chance alone a taster is right one time in three. If a panel identifies the odd sample significantly more often than one time in three, the samples are perceptibly different.

The number of correct answers needed depends on the number of tasters. At the conventional five percent significance level, a panel of 12 needs at least 8 correct identifications; 18 needs 10; 24 needs 13; 30 needs 15; 36 needs 18. These follow from the binomial distribution with a one-in-three chance of success.

\section{The limit of a null result}

A triangle test that does not reach significance does not show that samples are the same. It shows only that the panel did not detect a difference. A small panel will miss real but modest differences. Claims of similarity, as opposed to claims of difference, require larger panels and a stated maximum proportion of tasters who could distinguish the samples without the test detecting it.\footnote{Sensory science handles this with similarity testing, which sets the panel size in advance from the size of difference one wishes to rule out. The base-swap experiment should be designed this way if it is to support Hypothesis~\ref{h:late}.}

This matters for Hypothesis~\ref{h:late}. The hypothesis predicts that tasters cannot tell seasoned bases apart, which is a claim of similarity. A small panel that fails to distinguish them provides weak support at best.

\section{Sizing a similarity test}\label{sec:similarity}

A similarity test is designed around the largest difference one is willing to miss. In the usual model for the triangle test, a proportion $p_d$ of tasters can truly distinguish the samples and always choose correctly, while the rest guess and are right one time in three. The expected proportion of correct answers is then
\[
p_c = \tfrac{1}{3} + \tfrac{2}{3}\,p_d .
\]
To conclude that no more than a stated proportion $p_{\max}$ of tasters can distinguish the samples, the panel size $n$ and a cutoff $x$ are chosen so that, if $p_d$ were as large as $p_{\max}$, the chance of observing $x$ or fewer correct answers would be at most $\beta$. Similarity is concluded when the observed number correct is at most $x$. A second condition keeps the test from being uselessly strict: if the samples are truly indistinguishable ($p_d = 0$), the chance of observing $x$ or fewer correct should be at least $1-\alpha$.

Computing these from the binomial distribution with $\beta = 0.05$ and $\alpha = 0.20$ gives:

\begin{center}
\begin{tabular}{@{}ccc@{}}
\toprule
Largest $p_{\max}$ to rule out & Panel size $n$ & Conclude similarity if correct $\leq$ \\
\midrule
0.5 & 16 & 7 \\
0.4 & 25 & 10 \\
0.3 & 39 & 15 \\
0.2 & 86 & 32 \\
\bottomrule
\end{tabular}
\end{center}

Ruling out small proportions of distinguishers requires large panels. Experiment E1 runs three pairwise tests. If a single joint claim is to be made across all three, the risk $\beta$ is divided among them; at $\beta = 0.05/3$ per test, ruling out $p_{\max} = 0.3$ requires 58 tasters with a cutoff of 22, and ruling out $p_{\max} = 0.4$ requires 33 with a cutoff of 13.\footnote{These values follow directly from the binomial calculation described. Published tables for similarity testing use the same model and should agree closely; small differences can arise from how the two risks are balanced.}

A modest first run is realistic. A panel of about 25 to 40 tasters can rule out that as many as 30 to 40 percent of people distinguish the bases, which is enough to show whether the hypothesis is worth a larger study.

\section{Descriptive panels}

Descriptive analysis uses a trained panel to rate samples on agreed attributes, such as firmness, fibrousness, juiciness, and specific aroma notes, on numerical scales. It shows how samples differ, not just whether they do. For strand blocks, attributes include ease of separating along the grain, resistance across it, juiciness, and whether the texture changes during chewing.

\section{Acceptance}

Acceptance testing asks untrained consumers how much they like a product. It is the last test to run, because liking depends on everything at once and cannot guide structural development.

\section{Coding and order}

All samples are coded with random numbers, served in balanced order, and prepared identically except for the variable under test. Tasters rinse between samples. Where the base-swap test is run, the person serving should not know which base is which.

%----------------------------------------------------------------------
\chapter{The Limits of Equivalence}\label{ch:equivalence}

Status: an argument about evidence.

\section{What each test certifies}

No single test shows that a new food is equivalent to a reference food. Each test certifies similarity in one respect:

\begin{description}[leftmargin=1.5em]
\item[Proximate composition] certifies similar gross amounts of protein, fat, carbohydrate, and water.
\item[Centrifugal profile] certifies similar gross separation behavior after structure is destroyed.
\item[Texture profile analysis and shear force] certify similar mechanical response under a particular test.
\item[Cross-section imaging] certifies similar visible architecture.
\item[Triangle testing] certifies that a panel did not detect a difference, with the limits discussed in Chapter~\ref{ch:sensory}.
\item[Nutritional studies] certify similar digestion or health outcomes under particular conditions.
\end{description}

\section{The conjunction problem}

Two foods can match on every one of these and still differ in ways none of them captures, such as behavior when reheated, change during storage, or response to a different cooking method. Conversely, a food can differ from a reference on several measures and still serve the same purpose in a meal. Equivalence is therefore not a single property but a claim relative to a purpose, and the evidence required depends on the purpose.

\section{Stating claims properly}

The book adopts a rule for its own claims. A result is reported as similarity in a named respect, measured by a named method, under named conditions. The phrase ``equivalent to'' is not used without that qualification. The concept drawing that compares a centrifuged reference with a centrifuged prototype states the point in one line: a centrifuge reveals separation, not full equivalence.\footnote{The rule applies with particular force to comparisons with animal foods, where loose claims of equivalence are common in marketing and are the source of justified skepticism.}

%----------------------------------------------------------------------
\chapter{Nutrition and the Food Matrix}\label{ch:nutrition}

Status: raises questions to be tested; makes no nutritional claims.

\section{The food matrix}

The nutritional effect of a food depends on its structure as well as its composition. Nutrients enclosed within intact plant cell walls can be less accessible to digestion than the same nutrients released by milling or pressing. In almonds, for example, much of the lipid remains enclosed in cells that pass through chewing intact, and this cell-wall encapsulation limits how much of the fat is released during digestion~\cite{ellis2004}. The physical arrangement of nutrients, called the food matrix, has been reviewed as a determinant of how foods behave in processing and digestion and of their nutritional effects~\cite{aguilera2019}.

\section{The objection}

Fractionation destroys the native matrix of an ingredient, and reassembly builds a new one. This places the platform squarely within the debate over ultra-processed foods, which are commonly defined partly by their reliance on fractionated ingredients and industrial processes and have been associated in observational studies with poorer health outcomes. The mechanisms behind those associations are disputed, and the category is defined by processing and formulation rather than by structure. But the objection is serious and has to be answered by measurement, not by assertion.

A controlled trial gives the concern concrete form. In an inpatient randomized crossover study, Hall and colleagues fed adults diets of ultra-processed or unprocessed foods for two weeks each, with the diets matched for presented calories, sugar, fat, sodium, fiber, and macronutrients, and allowed participants to eat as much as they wanted. Participants ate more and gained weight on the ultra-processed diet~\cite{hall2019}. The study did not isolate fractionation, or any single property of processing, as the cause, and it does not show that every processed food has the same effect. What it shows is that matching the ingredients or nutrient content of a diet does not guarantee matching eating behavior. A reconstructed product therefore cannot be assumed nutritionally equivalent to a reference because its composition matches; how it is eaten has to be measured.

\section{A structural response}

The strand method differs from many fractionated foods in one relevant way: it builds a new matrix deliberately, with defined compartments and walls. Whether a built matrix slows digestion or increases satiety the way a native one does is an open question, but it is a question that can be asked of a strand block and not of a homogeneous paste. The structural variables of this book, $\phi$, $d$, $f$, may affect digestion rate as well as texture.

\section{Measured targets}

The following are proposed as measured properties of any product the platform makes, rather than assumed:

\begin{itemize}
\item composition, including fiber and sodium;
\item protein quality, from the amino acid profile of the protein sources used and their combination;
\item rate of starch and protein breakdown in standard laboratory digestion models;
\item for products that reach consumer testing, reported satiety compared with reference foods.
\end{itemize}

\section{Additives and sodium}

Part of the ultra-processed food concern concerns additives and high salt, sugar, and fat. The platform's reference formulations use a small number of familiar ingredients. Sodium is the main concern for savory products, since late-bound flavor relies partly on salt; sodium content is therefore recorded for every product and kept within the range of comparable home-cooked foods.\footnote{A product whose appeal depends on salt levels well above those of the dish it replaces would fail this chapter's standard regardless of its structure.}

%======================================================================
\part{Operations}
%======================================================================

\chapter{Batch Records}\label{ch:batch}

Status: proposed practice, following the process framework of Industrial Ecologies.

\section{The record as the product's history}

Every batch that leaves the platform is accompanied by a record of how it was made. The record serves three purposes: it allows a good batch to be repeated, it allows a bad batch to be traced to its cause, and it supports the release decision.

\section{Contents}

A batch record contains:

\begin{itemize}
\item the product route and the version of each formulation used;
\item lot identifiers for every input stream (Chapter~\ref{ch:intake});
\item for each operator, its settings and measured conditions: temperatures and times, pump rates, draw speeds, coiling parameters, press conditions;
\item for fermentations, culture identity and source, inoculation rate, and temperature and pH over time;
\item in-process measurements: preform sections, strand diameters, bend and peel tests where applicable;
\item finished product measurements: cross-sections, texture, cooking loss, sensory check;
\item cleaning and changeover records for shared equipment used before the batch;
\item the release decision, who made it, and on what basis.
\end{itemize}

\section{Setpoints, measurements, and calculated results}

The division introduced for fermentation in Chapter~\ref{ch:transformation} applies to every operator on the platform. For each parameter the record keeps, where applicable, three entries:

\begin{center}
\small
\begin{tabularx}{\linewidth}{@{}lXXX@{}}
\toprule
Operator & Setpoint & Measured & Calculated \\
\midrule
Fermentation & Incubator temperature & Vessel temperature and pH over time & Acidification rate, yield \\
Preform & Pump ratio $Q_c/(Q_c+Q_s)$ & Core and shell radii from sections & $\phi$ of the preform \\
Drawing & Roller speed ratio & Strand diameter at intervals & $\lambda$; $\phi$ of the strand \\
Coiling & Traverse and spindle speeds & Pitch on the finished coil & Helix angle $\theta$ \\
Pressing & Platen pressure and temperature & Block internal temperature over time & Fusion index $f$ from peel tests \\
Evaluation & Cooking procedure & Weights, forces, images & Cooking loss, $\mathcal{A}$, void fraction \\
\bottomrule
\end{tabularx}
\end{center}

The strand process makes the distinction unusually concrete. The core fraction set by pump rates is a setpoint; the core fraction read from a cross-section is a measurement; Hypothesis~\ref{h:draw} is precisely the claim that the two agree after drawing. A record that stored only one number for $\phi$ could not test the hypothesis at all.

\section{Standards and vocabularies}

Records are most useful when their terms mean the same thing from batch to batch and from one group to another. Several existing vocabularies cover parts of the platform. FoodOn, a harmonized food ontology, provides terms for foods, ingredients, and food processes intended to support traceability and data integration~\cite{foodon2018}. Work on the semantics of fermented dairy foods covers part of traditional fermentation~\cite{vitali2022}. For precision fermentation, the PREFER ontology provides an open vocabulary covering fermentation materials, processes, control settings, measurements, and calculated results~\cite{prefer2026}. PREFER's authors note that the earlier fermentation vocabularies were built for traditional and fermented-food applications and do not capture the specifics of precision fermentation, which is why they built a new one; the division matters here because the platform uses both kinds.

PREFER is constructed to conform to the Basic Formal Ontology, reuses terms from other community ontologies, and is maintained openly on a public repository to which others can propose terms and definitions. Its authors describe linked biofoundry datasets, knowledge graphs, and digital twins as future directions rather than completed work. It supports the case for explicit, reusable process records; it does not bear on whether the foods proposed here can be made.

No vocabulary known to this book covers the operations specific to the strand method: preform geometry, drawing, coiling, consolidation, and fusion. Nor do the fermentation vocabularies cover allergen cross-contact on shared lines or sensory measurements. Chapter~\ref{ch:vocabulary} proposes how the book's notation could be connected to the existing vocabularies rather than standing apart from them.

\section{Release}

A batch is released only when its finished measurements fall within the specification for its product and its record is complete. A batch that falls outside specification is held, investigated, and either reworked, diverted to another use, or discarded, and the decision is recorded.

\section{Records as research data}

On a platform that is also a research program, batch records are experimental data. Each batch is a point in the parameter space of Part V, with measured outcomes. Keeping records in a consistent, machine-readable form allows the product map to be built from routine production as well as from dedicated experiments.\footnote{A shared spreadsheet or database with one row per batch and one column per parameter and measurement is sufficient at bench and pilot scale.}

%----------------------------------------------------------------------
\chapter{Food Safety}\label{ch:safety}

Status: established principles applied to proposed processes. Any product served to others must be prepared under the food safety rules of the relevant jurisdiction; this chapter does not replace them.

\section{Hazard analysis per operator}

Food safety systems based on hazard analysis identify, for each step, the biological, chemical, and physical hazards that could occur and the controls that prevent them. On a shared platform the analysis is done per operator, since operators are shared across products, and then checked per route.

\section{Plant materials are not low risk by default}

Plant-based foods can carry pathogens from soil and water, and moist, protein-rich pastes support microbial growth as readily as meat does. Cooked rice, sprouted seeds, and cut produce are well-known sources of foodborne illness. Pastes, emulsions, and strands in this book should be treated with the same care as raw meat preparations.

\section{Cooling and holding}

Cooked foods must be cooled quickly through the temperature range in which bacteria grow fastest and held either hot or cold, not warm. As one example of a specific rule, the 2026 edition of the United States FDA Food Code, in \S3-501.14, requires cooked time/temperature control for safety food to be cooled from $57^\circ$C ($135^\circ$F) to $21^\circ$C ($70^\circ$F) within two hours, and to $5^\circ$C ($41^\circ$F) or below within six hours in total~\cite{foodcode2026}. The rule applies to foods in that category, which includes moist protein-rich preparations of the kind described in this book, not to every cooked food. The Food Code is a model code that jurisdictions choose whether and how to adopt, not a universal rule; the code in force where food is prepared is the one to follow.

\section{Enclosed cores}

The strand method creates a particular hazard. A core sealed inside a shell is shielded from air. A moist, low-acid core, such as a paste or an emulsion without added acid, held at room temperature in an oxygen-poor environment, could in principle allow the growth of \emph{Clostridium botulinum}, the organism responsible for botulism. The same concern applies to garlic or herbs stored in oil. Strand blocks should therefore be cooked promptly or kept refrigerated or frozen, and any product intended for storage at room temperature requires specific validation of acidity, water activity, or processing.

\section{Allergens and cross-contact}

A shared platform handles many ingredients on the same equipment. Priority allergens vary by jurisdiction but commonly include wheat and other gluten sources, soy, peanuts, tree nuts, sesame, milk, eggs, mustard, and others. The reference formulations use several of these. Each product's allergen content is declared, equipment is cleaned and verified between products with different allergen profiles, and schedules run allergen-free products before allergen-containing ones where possible.

\section{Wild and uncertified inputs}

Inputs such as cattail rhizomes gathered from the wild carry additional hazards: misidentification, since young cattail shoots can be confused with toxic plants such as some irises; contamination from polluted water, since aquatic plants can accumulate metals and other pollutants; and microbial contamination. Wild inputs should be sourced from known clean sites, positively identified, and tested where contamination is plausible.

\section{Physical hazards}

Mechanical equipment such as rollers, mandrels, and presses must not shed fragments into food. Food-contact parts are inspected, and blades used for sectioning are kept separate from production.

%----------------------------------------------------------------------
\chapter{By-Products and Stream Closure}\label{ch:closure}

Status: proposed.

\section{The principle}

Every fractionation produces streams the target product does not use. A platform is closed when every stream has a planned destination. An unplanned stream is a cost: it must be stored, disposed of, or wasted.

\section{Common by-products and uses}

\begin{table}[h]
\centering
\begin{tabularx}{\linewidth}{@{}lX@{}}
\toprule
By-product & Candidate uses \\
\midrule
Okara and legume pulp & Baked goods, fermentation substrate, strand shell filler, animal feed \\
Bran and hulls & Fiber enrichment of carriers, feed, compost \\
Fruit and vegetable pomace & Fiber, pectin extraction, fermentation, reconstructed fruit shells \\
Whey-like liquids from curd making & Fermentation media, cooking liquid, feed \\
Oilseed press cake & Protein flour, feed \\
Starch-rich sediment & Thickener, noodle and flatbread dough, sugar source by enzymatic treatment \\
Trimmings from strand blocks & Reforming into paste cores, sausages, fillings \\
\bottomrule
\end{tabularx}
\caption{By-products and candidate destinations.}
\end{table}

\section{Internal loops}

Some by-products return to the platform as inputs. Strand trimmings can be minced and used as paste cores. Starch sediment from protein extraction can become the carrier dough for flatbreads. Pomace can supply fiber and pectin for shells. Each internal loop reduces waste and raises the reuse of existing operators.\footnote{Internal loops also propagate problems: a contaminated by-product returned to production affects every product that uses it. Loops are recorded in the batch records like any other input.}

\section{External outlets}

Streams with no internal use go to feed, compost, or anaerobic digestion. These outlets have low value but prevent waste. A regional platform's economics depend in part on having reliable outlets nearby.

%----------------------------------------------------------------------
\chapter{Economics and Scale}\label{ch:economics}

Status: an analytical framework without cost data.

\section{Cost per operator}

The cost of a product is the sum over its route of the cost of each operator: equipment depreciation, energy, labor, cleaning, and losses, plus the cost of inputs. On a shared platform, equipment costs are spread over every product that uses the equipment. This is the economic content of the operator graph in Chapter~\ref{ch:graph}: heavily reused operators have low cost per product; operators used by one route carry their whole cost on that route.

\section{Utilization}

Equipment earns its cost only while it is working. A shared platform raises utilization by running several products through the same equipment, at the cost of changeovers, cleaning, and scheduling. The balance between utilization gained and changeover time lost determines whether sharing pays.

\section{Scale}

Food manufacturing has strong economies of scale. Large plants spread fixed costs over enormous volumes, buy inputs cheaply, and run continuous lines without changeovers. A small regional platform cannot match their unit costs for commodity products. Its possible advantages lie elsewhere: flexibility to use local inputs, lower transport, the ability to make products too varied or specialized for large lines, and the ability to learn quickly by changing routes.

\section{A break-even condition}

Without cost data the chapter cannot say whether the strand method pays, but it can say what would have to be true. Let each operator $k$ have a fixed cost $F_k$ per year and a variable cost $v_k$ per kilogram processed, and let it process $Q_k$ kilograms per year across all routes that use it. A product whose route passes through operators $k \in \mathcal{R}$, with input cost $m$ per kilogram and overall yield $y$, costs approximately
\[
c = \frac{m}{y} + \sum_{k \in \mathcal{R}} \left( \frac{F_k}{Q_k} + v_k \right)
\]
per kilogram of product. Shared operators have large $Q_k$ and therefore small fixed cost per kilogram; the coiling station and press, used only by the strand route, carry their fixed cost against strand output alone.

If $c_s$ is the cost of a strand block and $c_t$ the cost of a comparable dish base made from textured vegetable protein, the strand block is worth making only if the premium buyers will pay for its texture, $\Delta p$, satisfies
\[
\Delta p \geq c_s - c_t .
\]
The data needed to evaluate this are the operator costs and throughputs from the pilot line and a measured willingness to pay, which requires acceptance testing of a product that already works. The condition is stated here so that those measurements are collected in a form that answers it.

\section{Why a better process can lose}

Textured vegetable protein is cheap, shelf-stable, and adequate for many dishes. A strand block may have better structure but will cost more per kilogram, at least initially, because it involves more steps. Products compete on price, convenience, familiarity, and shelf life as well as on quality. A process that gives a better result can still lose to one that gives an adequate result at lower cost, and many technically superior food processes have. The strand method should be evaluated against this standard: not only whether it makes better structure, but whether the improvement is worth its cost to the people who would buy it.\footnote{The same standard applies to the platform as a whole. Flexibility has value only if it is used.}

%----------------------------------------------------------------------
\chapter{Availability Before Restriction}\label{ch:availability}

Status: a normative argument, drawn from Debt Before Deflation.

\section{The argument}

Proposals to reduce the environmental and ethical costs of established food systems often take the form of restrictions: taxes, bans, quotas, or regulations that raise the price or reduce the supply of existing foods. The argument of this chapter is that such restrictions are fair only when acceptable substitutes are genuinely available: affordable, accessible, familiar enough to be used, and satisfying enough to be chosen.

Where substitutes are not available, restrictions fall hardest on those with the least ability to absorb higher prices or to change their diets, and they tend to provoke resistance that sets back the goals they were meant to serve.

\section{Texture and flavor as technical problems}

The argument connects to the rest of the book through a factual claim: that much of what makes established foods hard to replace is texture and flavor, and that texture and flavor are partly problems of processing and chemistry. If that is so, then developing processes that produce satisfying textures from widely available ingredients is part of making substitutes available, and therefore part of the preconditions for any fair change.

\section{The limits of the argument}

The argument has limits that should be stated. It does not say which restrictions, if any, are justified; that depends on judgments about environmental, health, and ethical costs on which people reasonably disagree. It does not say that substitutes must be identical to what they replace. And it does not claim that technical availability is sufficient: cultural attachment, habit, and trust matter, and a food can be available without being accepted.\footnote{Others argue that restrictions themselves drive the innovation that produces substitutes, so that waiting for availability delays both. The disagreement is partly empirical and is noted here rather than resolved.}

The claim is narrower: whatever changes are pursued, work that makes good substitutes available widens the range of options and reduces the burden of change on those least able to bear it. This book is a contribution to that work.

%======================================================================
\part{Research Program}
%======================================================================

\chapter{Practiced, Proposed, Hypothesized}\label{ch:ledger}

Status: a ledger of the book's claims.

\begin{table}[h]
\centering
\small
\begin{tabularx}{\linewidth}{@{}Xll@{}}
\toprule
Item & Category & Where tested \\
\midrule
Fried protein base in broth or oil with onion, herbs, spices & Practiced & Appendix B \\
Carrier, matrix, inclusion as a working description & Practiced & (descriptive) \\
Late binding of perceived identity & Hypothesized (H1) & E1 \\
Phase fraction as an independent control variable & Hypothesized (H2) & E7 \\
Emulsion analogy, milk to mayonnaise & Analogy & E2 (calibration) \\
Centrifugal intake profiles & Proposed & Routine use \\
Fruit milking as defined & Proposed & Product trials \\
Rolling hotplate compression & Proposed & Product trials \\
Shared operator platform & Proposed & Operation \\
Coaxial preform construction & Proposed & E3, E4 \\
Shell setting routes & Proposed & E3 \\
Draw invariance of $\phi$ & Hypothesized (H3) & E4 \\
Coiling at controlled $\theta$ & Proposed & E5 \\
Shell fusion surviving cooking & Hypothesized (H4) & E6 \\
Product map regions & Hypothesized & E7 \\
Reconstructed fruit & Proposed & Later trials \\
Nutritional behavior of built matrices & Open question & Later studies \\
\bottomrule
\end{tabularx}
\caption{Status of every principal claim.}
\end{table}

The ledger is to be updated as results arrive. An item moves from hypothesized to supported, qualified, or rejected only on the basis of recorded experiments, and the change is dated.

%----------------------------------------------------------------------
\chapter{Experiment Sequence}\label{ch:sequence}

Status: proposed.

The sequence is ordered so that each stage fixes some variables before the next varies others, and so that the cheapest experiments come first. Each stage has a gate: a result that must be obtained before later stages are worth running.

\section*{E1. Base swap}

\textbf{Question.} Can tasters distinguish textured vegetable protein, lentils, and bread prepared identically with a complex seasoning? \textbf{Method.} Prepare each base by the protocol of Appendix B with identical oil, onion, herbs, and spice mixture. Run triangle tests for each pair, with a panel sized for similarity testing as in Section~\ref{sec:similarity}; a first run of 25 to 40 tasters, with the proportion of distinguishers to be ruled out stated in advance. Repeat with a simple seasoning. \textbf{Measurements.} Correct identifications; tasters' descriptions of differences. \textbf{Gate.} None; E1 is independent of the strand work. Its result bears on Hypothesis~\ref{h:late} and on how much of product identity can be left to finishing.

\section*{E2. Emulsion series}

\textbf{Question.} How does consistency change with $\psi$ at fixed emulsifier and mixing? \textbf{Method.} Prepare oil-in-water emulsions at stepped oil fractions. \textbf{Measurements.} Flow behavior, shape holding, stability over a day and after heating. \textbf{Gate.} Identifies usable emulsion cores for later stages.

\section*{E3. Shell setting}

\textbf{Question.} Which shell and setting route satisfies containment, flexibility, and bondability together? \textbf{Method.} Form single strands from each pairing in Chapter~\ref{ch:formulations}. \textbf{Measurements.} Leakage; minimum bend radius; peel force after pressing strand pairs. \textbf{Gate.} At least one shell must pass all three. If none does, the process stops here and the shell problem becomes the research program.

\section*{E4. Draw}

\textbf{Question.} Does drawing preserve $\phi$? \textbf{Method.} Draw preforms of the passing shell at several $\lambda$. \textbf{Measurements.} $\phi$ along the strand from cross-sections; breakage; core beading. \textbf{Gate.} A draw ratio range over which $\phi$ holds within a stated tolerance.

\section*{E5. Coiling}

\textbf{Question.} Can strands be coiled at controlled $\theta$ without damage? \textbf{Method.} Coil at near-hoop and crossed angles on mandrels sized from the bend test. \textbf{Measurements.} Shell cracking; pitch uniformity; handling. \textbf{Gate.} Reliable coiling at two angles.

\section*{E6. Consolidation}

\textbf{Question.} Do shells fuse while cores remain distinct, and does this survive cooking? \textbf{Method.} Press coiled blocks under varied $P$, $T$, $t$, with and without oiled interfaces and binders. \textbf{Measurements.} Void fraction, wall continuity, core integrity, raw and cooked; fusion index from peel tests. \textbf{Gate.} At least two fusion levels reproducibly achieved, with structure surviving cooking. This stage tests Hypothesis~\ref{h:fusion}.

\section*{E7. Texture across the map}

\textbf{Question.} Do $\phi$, $\theta$, and $f$ change cross-sectional structure and directional mechanical response? \textbf{Method.} Make blocks across the trial matrix of Chapter~\ref{ch:params} with other variables fixed. \textbf{Measurements.} Anisotropy index $\mathcal{A}$, texture profiles in two directions, cooking loss, cross-sections, triangle tests between neighboring conditions. \textbf{Outcome.} Measured boundaries for the product map, and a test of Hypothesis~\ref{h:phase}.

\section*{After E7}

Only after the structural work would descriptive sensory panels, nutritional measurements, and comparisons with reference foods be worth running. Perceptual identity and nutrition are separate questions and are not advanced by the structural results alone.

%----------------------------------------------------------------------
\chapter{A Vocabulary for Structured Foods}\label{ch:vocabulary}

Status: proposed.

\section{Why a vocabulary}

A research program that other people can contribute to needs terms that mean the same thing to everyone who uses them. The notation of this book, $\phi$, $\lambda$, $\theta$, $f$, $\mathcal{A}$, and the operator names of Part III, is a start. This chapter describes how it could be made into a shared vocabulary that connects to existing ones.

\section{Domain and application layers}

The PREFER ontology describes a layered arrangement~\cite{prefer2026}. A domain ontology defines core concepts once. Individual laboratories keep their own local terms and connect them to the domain ontology through an application ontology, a set of logical mappings. One of its examples is a laboratory that records a fermentation only as run at ``default conditions.'' The application ontology expands that phrase into explicit temperature and pH setpoints with units and values, so the laboratory's records can be compared with those of laboratories that record setpoints separately.

The same arrangement suits this platform. Existing domain vocabularies cover foods and ingredients (FoodOn~\cite{foodon2018}) and fermentation (PREFER and its predecessors). A small application vocabulary would define the terms specific to the strand method and to the platform's records, and map them to the domain vocabularies wherever they overlap. The mapping would do for the strand process what PREFER's example does for fermentation: a record stating that a block was ``pressed as usual'' would be expanded into explicit setpoints for $P$, $T$, and $t$, and a record stating a strand was ``half core'' into a setpoint for $\phi$ with its measured counterpart.

\section{Relations between terms}

Most vocabularies are organized by three relations: one term is broader than another, narrower than another, or the same as another. Placing the book's terms in existing vocabularies means deciding these relations: whether \emph{consolidation} is narrower than an existing term for pressing, whether \emph{fusion index} is the same as any existing bond-strength measure, whether \emph{core--shell strand} is narrower than an existing term for filled or co-extruded products.

Language models have been evaluated on proposing such relations. In a benchmark built from the IEEE Thesaurus, several models identified broader, narrower, and same-as relations between pairs of engineering topics with high accuracy, and smaller models guided by careful prompting performed well at low computational cost~\cite{aggarwal2026}. That result concerns research topics in engineering, not food process terms, and a benchmark score does not guarantee correct individual decisions. Its use here would be limited to proposing candidate relations for review. Every mapping entered in the vocabulary would be checked by a person who understands both terms.

\section{Building a literature base}

Much relevant prior work, on extrusion conditions, protein gels, alginate setting, fiber formation, and texture measurement, exists only as narrative papers and tables. Compiling it into records with the same fields as the platform's batch records would let prior results be placed on the product map alongside new ones.

Recent work on extracting structured records from scientific papers with language models is relevant to how this might be done. One study in yeast metabolic engineering combined targeting of relevant sections, expansion of context, an adversarial re-checking step that asks whether the source text explicitly supports each extracted value, and validation of records against an ontology. Against expert-curated records it reported large reductions in mismatched and hallucinated entries compared with simpler extraction~\cite{yang2027}. Its authors are explicit about the limits: a single domain, a small verified subset, conformance to constraints that does not establish correctness against the source, and a continuing need for human review where errors matter.

Those limits set the rule for any literature base built for this book. An extracted record is a candidate until a person has checked each value against the source. A record that conforms to the vocabulary's structure is not thereby correct. Extracted literature records are kept distinct from the platform's own measured records, and the product map shows which points come from which.\footnote{The same distinction between setpoint, measurement, and calculation applies to literature records: a paper's stated process conditions are setpoints unless the paper reports what was measured.}

\section{Open development}

A vocabulary that others can use has to be open to their corrections. PREFER is maintained on a public repository with an issue process for proposing and revising terms~\cite{prefer2026}. The application vocabulary for structured foods would be published the same way, beginning with the notation and glossary of this book, so that terms can be proposed, disputed, and refined by anyone working on the problems of Part VIII.

%----------------------------------------------------------------------
\chapter{Open Problems}\label{ch:open}

Status: unresolved.

\paragraph{Shell chemistry.} No shell formulation is yet known to satisfy containment, flexibility, and bondability together. This is the problem most likely to decide whether the strand method works at all.

\paragraph{Fusion after setting.} Setting and fusion pull in opposite directions. Sequences that separate them, such as drying and remoistening, or ionic skins over heat-fusing protein, are promising but untested.

\paragraph{Continuous drawing.} Hand drawing is slow and variable. Continuous drawing at useful speed without necking or breakage is an open engineering problem, likely limited by setting time.

\paragraph{Core leakage during cooking.} Fat and emulsion cores melt or thin when heated. Whether thin shells contain them through cooking determines whether the high-$\phi$ regions of the product map are reachable.

\paragraph{Shelf life and safety of enclosed cores.} Moist cores sealed in shells require refrigeration or validated preservation. Room-temperature storage is out of reach without further work.

\paragraph{Gluten-free shells.} The shell best suited to drawing depends on gluten. A shell without it that draws as well would widen the method's use.

\paragraph{Scale of structure.} Strands of a millimeter or more reproduce bundle-scale structure. Whether finer strands can be made and assembled at useful rates is unknown.

\paragraph{The transitional food map.} The stable and unstable regions between familiar dishes have not been mapped, and the method for doing so is only outlined in Chapter~\ref{ch:transitional}.

\paragraph{A shared vocabulary.} Terms specific to strand forming and to shared-platform records have no established vocabulary. Building one that maps onto existing food and fermentation ontologies is outlined in Chapter~\ref{ch:vocabulary}.

\paragraph{Built matrices and nutrition.} Whether deliberately built compartments affect digestion and satiety as native food structure does is an open scientific question.

Each of these is stated so that it can be worked on independently. Progress on any one advances the program, whether or not the others are solved.

%======================================================================
\appendix
%======================================================================

\chapter{Parameter Tables}\label{app:params}

\section{Structural and process parameters}

\begin{center}
\small
\begin{tabularx}{\linewidth}{@{}llllX@{}}
\toprule
Symbol & Quantity & Unit & Kind & How obtained \\
\midrule
$R_0, r_0$ & Preform outer and core radius & mm & Set / measured & Set by tooling or pump ratio; measured from preform sections \\
$\phi$ & Core area fraction & -- & Calculated & $(r/R)^2$ from sections, or $Q_c/(Q_c+Q_s)$ as setpoint \\
$\lambda$ & Diameter draw ratio & -- & Calculated & Preform diameter over strand diameter; $\sqrt{v_2/v_1}$ as setpoint \\
$d$ & Strand diameter & mm & Measured & Calipers at stated intervals \\
$\varepsilon_{\max}$ & Maximum surface strain & -- & Calculated & $R/\rho$ at smallest dowel survived in bend test \\
$\rho$ & Mandrel radius & mm & Set & Tooling \\
$\theta$ & Helix angle & degrees & Set / calculated & Setpoint from traverse and spindle speeds; checked from measured pitch \\
$p$ & Pitch & mm & Measured & Axial advance per turn on finished coil \\
$n$ & Layer count & -- & Set & Counted \\
$\eta$ & Packing fraction before pressing & -- & Measured & Strand area over total area in unpressed sections \\
$P$ & Press pressure & kPa & Set / measured & Load over platen area \\
$T$ & Press temperature & $^\circ$C & Set / measured & Platen setpoint; block probe log \\
$t$ & Press dwell & min & Set & Timer \\
$f$ & Fusion index & -- & Calculated & $F_{\text{sep}}/F_{\text{ref}}$ under a stated test geometry, with failure location (Protocol~\ref{prot:peel}) \\
$\psi$ & Emulsion dispersed fraction & -- & Set & Oil volume over total volume \\
\bottomrule
\end{tabularx}
\end{center}

\section{Outcome measures}

\begin{center}
\small
\begin{tabularx}{\linewidth}{@{}lllX@{}}
\toprule
Symbol & Quantity & Unit & Method \\
\midrule
$\mathcal{A}$ & Anisotropy index & -- & Shear force parallel over perpendicular (Chapter~\ref{ch:texture}) \\
-- & Hardness, cohesiveness, springiness, chewiness & N, -- & Texture profile analysis \\
-- & Cooking loss & \% & Weight before and after standard cooking \\
-- & Void fraction & \% & Void area in pressed sections \\
-- & Core integrity & score & Sections scored 0 (cores lost) to 3 (cores intact) \\
-- & Wall continuity & score & Sections scored 0 (all boundaries visible) to 3 (none visible) \\
-- & Leakage & g & Core material recovered after standard pressure on single strands \\
-- & Triangle test result & correct / $n$ & Chapter~\ref{ch:sensory} \\
\bottomrule
\end{tabularx}
\end{center}

\section{Measured ranges}

This section will hold the ranges over which each parameter has been achieved and the outcomes observed, filled from batch records as experiments proceed. It is deliberately empty in the first edition.

%----------------------------------------------------------------------
\chapter{Domestic Protocols}\label{app:protocols}

The protocols in this appendix record practiced methods in a form that can be repeated and varied. Quantities are given as starting values; the variables to record are listed with each protocol.

\begin{protocol}[Fried protein base]\label{prot:tvp}
\textbf{Purpose.} A savory, meat-like base for sauces, pasta, and fillings, and the reference preparation for Experiment E1.

\textbf{Ingredients.}
\begin{itemize}
\item textured vegetable protein, dry: 100\,g;
\item broth or seasoned water, hot: 150 to 200\,g;
\item neutral oil: 20 to 40\,g;
\item onion, finely chopped: 100\,g;
\item garlic, herbs, and spice mixture: to a recorded formula;
\item salt: to a recorded amount;
\item optional bitter or roasted addition, such as instant coffee or cocoa: 1 to 3\,g.
\end{itemize}

\textbf{Method.}
\begin{enumerate}
\item Pour the hot broth over the textured vegetable protein and let it absorb for about ten minutes, until no dry pieces remain. Press out any excess liquid and weigh the hydrated base.
\item Heat the oil in a wide pan over medium-high heat. Add the onion and cook until soft and beginning to brown.
\item Add the hydrated base in a single layer and leave it undisturbed until the underside browns, then turn and repeat.
\item Add garlic, spices, and any bitter addition, stirring briefly to release their aroma in the oil.
\item Season with salt and, if a sauce is wanted, deglaze with a little broth.
\end{enumerate}

\textbf{Variables to record.} Weights of all ingredients; broth composition; hydration time; pan temperature where measurable; browning time per side; final weight; spice formula, including the number of components.
\end{protocol}

\begin{protocol}[Lentil and bread variants]\label{prot:variants}
\textbf{Lentils.} Cook brown or green lentils until just tender and drain well. Use 250\,g of cooked lentils in place of the hydrated base in Protocol~\ref{prot:tvp}, frying until the surfaces crisp.

\textbf{Bread.} Tear day-old bread into small pieces and moisten lightly with broth. Use 150\,g in place of the hydrated base, frying until golden.

\textbf{Variables to record.} As for Protocol~\ref{prot:tvp}, with the base substituted.
\end{protocol}

\begin{protocol}[Base-swap preparation for E1]\label{prot:e1}
\textbf{Purpose.} To prepare samples that differ only in their base.
\begin{enumerate}
\item Prepare a single batch of onion, oil, and spice mixture sufficient for all three bases, and divide it by weight into three equal parts.
\item Prepare the three bases by Protocols~\ref{prot:tvp} and~\ref{prot:variants}, each with one part of the shared mixture, in the same pan type and with the same timings.
\item Run three pairwise triangle tests: textured vegetable protein against lentils, textured vegetable protein against bread, and lentils against bread. Each taster completes all three, in randomized order.
\item For each triangle, assign the presentation order at random from the six possible arrangements (for bases A and B: AAB, ABA, BAA, BBA, BAB, ABB), balanced so that each arrangement occurs equally often across the panel.
\item A coder, working away from the tasting area, portions each preparation into identical containers labeled only with random three-digit codes and keeps the code key. The server receives coded trays only and has no access to the key until all responses are collected.
\item Serve all samples at the same temperature and in the same quantity.
\end{enumerate}
\textbf{Simple-seasoning arm.} Repeat with only oil, onion, and salt.
\end{protocol}

\begin{protocol}[Centrifugal profile]\label{prot:centrifuge}
\begin{enumerate}
\item Weigh 20\,g of sample and 40\,g of water.
\item Blend for a fixed time at a fixed setting.
\item Fill graduated tubes to a fixed volume.
\item Spin at a fixed speed for a fixed time, recorded.
\item Read the volumes of the oil-rich, aqueous, and sediment layers.
\end{enumerate}
\textbf{Note.} The profile characterizes gross composition only (Chapter~\ref{ch:equivalence}).
\end{protocol}

\begin{protocol}[Bend test]\label{prot:bend}
\begin{enumerate}
\item Prepare dowels of graded diameter, for example 5, 10, 15, 20, 30, and 40\,mm.
\item Wrap a single strand once around each dowel, starting with the largest.
\item Inspect the outer surface for cracks under good light.
\item Record the smallest dowel survived without cracking, and compute $\varepsilon_{\max} = R/(r_{\text{dowel}} + R)$, since the strand's centerline lies one strand radius outside the dowel surface.
\end{enumerate}
\end{protocol}

\begin{protocol}[Separation test for the fusion index]\label{prot:peel}
\begin{enumerate}
\item \textbf{Pair specimen.} Lay two strands of equal length side by side and press them under the conditions being tested, keeping a short unbonded tail at one end of each, for example by placing a strip of baking paper between them there. Record the strand diameter, the bonded length $L_b$, and the tail length.
\item Clamp the two tails in opposing clamps, one fixed and one attached to a spring scale or force gauge, so that the strands are pulled apart from the tail end along the bond.
\item Pull at a steady, recorded rate and record the peak force $F_{\text{sep}}$.
\item Note where failure occurred: along the interface, within a shell, or through a core.
\item \textbf{Reference specimen.} Take a single strand of the same type, given the same press and heat treatment without a partner. Clamp its two ends in the same clamps with a recorded free length between them, and pull it to breaking at the same rate. Record the peak force $F_{\text{ref}}$ and where it broke; a break at the clamp is recorded as such and the test repeated.
\item Compute $f = F_{\text{sep}}/F_{\text{ref}}$. If the pair failed through a strand rather than at the interface, record $f \geq 1$ with cohesive failure.
\end{enumerate}
\textbf{Note.} The ratio depends on both specimen geometries and is used to compare conditions tested the same way, not to report an intrinsic bond strength. Report $f$ together with $L_b$, the reference free length, and the failure mode of each specimen.
\end{protocol}

%----------------------------------------------------------------------
\chapter{Concept Drawings}\label{app:drawings}

Five concept drawings and a cover illustration accompany the proposal. They are illustrations of ideas, not records of apparatus, and nothing shown in them has been built. Each carries its own disclaimer to that effect.

\begin{description}[leftmargin=1.5em,style=nextline]
\item[Cover.] A pilot facility in which two strand combs feed a coiling drum, with a schematic on a control screen showing the transition from an unstructured to a fibrous cross-section.
\item[Drawing 1: bilateral feeds.] Two hoppers feeding pastes from opposite sides; combs dividing each stream into many strands; the strands coiled together on a drum; and a finished block showing pull-apart texture in one region and sliceable resistance in another.
\item[Drawing 2: the heated passage.] Formulation, surface heating of an extruded strand, the resulting skin-and-core cross-section, and strands kept separate on a guide.
\item[Drawing 3: centrifugal comparison.] A reference food and a prototype blended and spun, with oil-rich, water-rich, and sediment-rich fractions compared. Its caption states that a centrifuge reveals separation, not full equivalence.
\item[Drawing 4: coil, press, season, test.] Strands wound from both sides onto a mandrel, pressed into a cohesive piece, seasoned on the surface and between accessible layers, and tested by pulling along the fibers and cutting across them.
\item[Drawing 5: seasoning families.] A neutral base divided among spice families, flavor delivered to surface and interlayers, and a coded blind comparison with a scorecard. Its caption states that spice is an experimental variable, not a substitute for texture.
\end{description}

% To include the drawings, place the image files alongside this source
% and uncomment the lines below. Add \usepackage{graphicx} to the preamble.
% \begin{figure}[p]\centering\includegraphics[width=\linewidth]{01-drawing.png}\caption{Drawing 1.}\end{figure}
% \begin{figure}[p]\centering\includegraphics[width=\linewidth]{02-drawing.png}\caption{Drawing 2.}\end{figure}
% \begin{figure}[p]\centering\includegraphics[width=\linewidth]{03-drawing.png}\caption{Drawing 3.}\end{figure}
% \begin{figure}[p]\centering\includegraphics[width=\linewidth]{04-drawing.png}\caption{Drawing 4.}\end{figure}
% \begin{figure}[p]\centering\includegraphics[width=\linewidth]{05-drawing.png}\caption{Drawing 5.}\end{figure}

The drawings anticipate several points that the text develops. Drawing 2 identifies shell setting as a separate step, the problem treated in Chapter~\ref{ch:shell}. Drawing 3 states the limit of compositional comparison argued in Chapter~\ref{ch:equivalence}. Drawing 4 separates testing along and across the fibers, the basis of the anisotropy index. Drawing 5 notes that surface seasoning reaches outer layers and interfaces but not the center, the problem that core seasoning addresses in Chapter~\ref{ch:finishing}. The drawings also show things the text qualifies: the finished block in Drawing 1 has no central bore, a problem discussed in Chapter~\ref{ch:coil}, and the drawings use animal-food names for textures that the text describes by structure.

%----------------------------------------------------------------------
\chapter{Glossary}\label{app:glossary}

\begin{description}[leftmargin=1.5em,style=nextline]
\item[Anisotropy index ($\mathcal{A}$)] Shear force along the strands divided by shear force across them.
\item[Binder] A material added to hold pieces together, such as bread in meatloaf or a protein slurry between strand layers.
\item[Calculated value] A value derived after the fact from measurements.
\item[Carrier] The starch-based body of a dish that supplies bulk and bite.
\item[Coiling] Laying strands onto a rotating mandrel along helical paths.
\item[Consolidation] Pressing coiled strands into a cohesive block.
\item[Core] The inner phase of a composite strand.
\item[Core fraction ($\phi$)] The proportion of a strand's cross-sectional area occupied by core.
\item[Cross-contact] The transfer of an allergen from one food to another through shared equipment or handling.
\item[Draw-down] Stretching a preform to reduce its diameter.
\item[Draw ratio ($\lambda$)] Preform diameter divided by strand diameter.
\item[Emulsion] A dispersion of droplets of one liquid in another with which it does not mix.
\item[Endomysium, perimysium, epimysium] The connective tissue sheaths around muscle fibers, fiber bundles, and whole muscles.
\item[Fascicle] A bundle of muscle fibers.
\item[Fruit milking] The production of a stable, pourable suspension or emulsion from plant tissue for use where milk would be used (Chapter~\ref{ch:fractionation}).
\item[Fusion index ($f$)] Peak force to separate two pressed strands divided by peak force to break one strand, measured in a stated geometry; an operational, dimensionless ratio.
\item[Helix angle ($\theta$)] The angle between a coiled strand and the mandrel axis.
\item[Hoop coiling] Coiling at a helix angle near $90^\circ$, so that strands run around the circumference.
\item[Inclusion] A discrete piece distributed through a dish's matrix.
\item[Late binding] The hypothesis that perceived identity is determined mainly at the finishing stage (Hypothesis~\ref{h:late}).
\item[Lineage] The path of a particular batch through the operator graph, with the lots and runs at each step.
\item[Maillard reaction] The network of reactions between sugars and amino acids on heating that produces browning and roasted flavor.
\item[Matrix] The sauce, gravy, or binding phase of a dish.
\item[Measured value] A value recorded from what actually happened during a process.
\item[Operator] A unit operation that transforms one or more streams in a known way.
\item[Operator graph] The network of operators and the product routes through them.
\item[Packing fraction ($\eta$)] The proportion of a cross-section filled by strands before pressing.
\item[Precision fermentation] The use of microorganisms, often engineered, to produce specific ingredients.
\item[Preform] A thick composite rope in which core and shell are arranged in their intended proportions before drawing.
\item[Retronasal aroma] Aroma reaching the nose from the mouth during chewing and swallowing, perceived as taste.
\item[Route] An ordered sequence of operators that makes a product.
\item[Setpoint] The value an operator is asked to achieve.
\item[Shell] The outer phase of a composite strand.
\item[Stream] A supply of one kind of input material, characterized at intake.
\item[Triangle test] A discrimination test in which tasters identify the odd sample among three.
\end{description}

%----------------------------------------------------------------------
\chapter{Recording Sheets}\label{app:sheets}

Each sheet corresponds to one experiment in Chapter~\ref{ch:sequence}. Every sheet begins with the same header.

\begin{center}
\begin{tabularx}{\linewidth}{@{}lX@{}}
\toprule
Date & \\
Operator & \\
Experiment and run number & \\
Formulation versions & \\
Input lot identifiers & \\
Equipment and last cleaning & \\
\bottomrule
\end{tabularx}
\end{center}

\section*{E1. Base swap}
\begin{center}
\small
\begin{tabularx}{\linewidth}{@{}lXXXl@{}}
\toprule
Taster & Pair tested & Codes served & Code chosen & Correct \\
\midrule
 & & & & \\ \midrule
 & & & & \\ \midrule
 & & & & \\ \midrule
 & & & & \\
\bottomrule
\end{tabularx}
\end{center}
Seasoning arm (complex or simple): \rule{4cm}{0.4pt}\qquad Total correct: \rule{2cm}{0.4pt} of \rule{1.5cm}{0.4pt}

\section*{E2. Emulsion series}
\begin{center}
\small
\begin{tabularx}{\linewidth}{@{}lXXXX@{}}
\toprule
$\psi$ & Emulsifier and amount & Flow or shape holding & Stable at 24\,h & Stable after heating \\
\midrule
 & & & & \\ \midrule
 & & & & \\ \midrule
 & & & & \\
\bottomrule
\end{tabularx}
\end{center}

\section*{E3. Shell setting}
\begin{center}
\small
\begin{tabularx}{\linewidth}{@{}llXXXX@{}}
\toprule
Shell & Core & Route and conditions & Leakage (g) & Smallest dowel (mm) & Peel force (N) \\
\midrule
 & & & & & \\ \midrule
 & & & & & \\ \midrule
 & & & & & \\
\bottomrule
\end{tabularx}
\end{center}

\section*{E4. Draw}
\begin{center}
\small
\begin{tabularx}{\linewidth}{@{}lXXXX@{}}
\toprule
Preform $\phi$ (setpoint) & $\lambda$ & $d$ at intervals (mm) & $\phi$ measured (mean, range) & Breaks, beading \\
\midrule
 & & & & \\ \midrule
 & & & & \\ \midrule
 & & & & \\
\bottomrule
\end{tabularx}
\end{center}

\section*{E5. Coiling}
\begin{center}
\small
\begin{tabularx}{\linewidth}{@{}lXXXX@{}}
\toprule
$\rho$ (mm) & $\theta$ setpoint & Pitch measured & $\theta$ calculated & Cracking, handling notes \\
\midrule
 & & & & \\ \midrule
 & & & & \\ \midrule
 & & & & \\
\bottomrule
\end{tabularx}
\end{center}

\section*{E6. Consolidation}
\begin{center}
\small
\begin{tabularx}{\linewidth}{@{}lXXXXX@{}}
\toprule
$P$, $T$, $t$ & Interface treatment & Void \% & Wall score & Core score (raw / cooked) & $f$ \\
\midrule
 & & & & & \\ \midrule
 & & & & & \\ \midrule
 & & & & & \\
\bottomrule
\end{tabularx}
\end{center}

\section*{E7. Texture across the map}
\begin{center}
\small
\begin{tabularx}{\linewidth}{@{}llllXXX@{}}
\toprule
$\phi$ & $d$ & $\theta$ & $f$ & $\mathcal{A}$ & Cooking loss (\%) & Notes \\
\midrule
 & & & & & & \\ \midrule
 & & & & & & \\ \midrule
 & & & & & & \\
\bottomrule
\end{tabularx}
\end{center}

\backmatter
\begin{thebibliography}{99}
\addcontentsline{toc}{chapter}{References}

\bibitem{aggarwal2026}
T.~Aggarwal, A.~Salatino, F.~Osborne, and E.~Motta.
\newblock Large language models for scholarly ontology generation: An extensive analysis in the engineering field.
\newblock \emph{Information Processing and Management}, 63:104262, 2026.
\newblock Available online 21 July 2025. doi:10.1016/j.ipm.2025.104262.

\bibitem{aguilera2019}
J.~M. Aguilera.
\newblock The food matrix: Implications in processing, nutrition and health.
\newblock \emph{Critical Reviews in Food Science and Nutrition}, 59(22), 2019.
\newblock doi:10.1080/10408398.2018.1502743.

\bibitem{prefer2026}
T.~Amigó, S.~Z.~K. Tan, A.~L. Phanthanourak, S.~Schulz, P.~D. Colaianni, D.~M. Maszczyk, E.~Milesi, I.~Schlembach, M.~Semenov Petrov, M.~Reventós Montané, L.~K. Nielsen, J.~Förster, B.~Ø. Palsson, S.~Sudarsan, and A.~Santos.
\newblock PREFER: An ontology for the PREcision FERmentation community.
\newblock arXiv:2602.16755v1, 2026.

\bibitem{bourne2002}
M.~C. Bourne.
\newblock \emph{Food Texture and Viscosity: Concept and Measurement}, 2nd edition.
\newblock Academic Press, 2002.

\bibitem{dekkers2018}
B.~L. Dekkers, R.~M. Boom, and A.~J. van~der Goot.
\newblock Structuring processes for meat analogues.
\newblock \emph{Trends in Food Science \& Technology}, 81:25--36, 2018.
\newblock doi:10.1016/j.tifs.2018.08.011.

\bibitem{foodon2018}
D.~M. Dooley et al.
\newblock FoodOn: A harmonized food ontology to increase global food traceability, quality control and data integration.
\newblock \emph{npj Science of Food}, 2:23, 2018.

\bibitem{ellis2004}
P.~R. Ellis, C.~W.~C. Kendall, Y.~Ren, C.~Parker, J.~F. Pacy, K.~W. Waldron, and D.~J.~A. Jenkins.
\newblock Role of cell walls in the bioaccessibility of lipids in almond seeds.
\newblock \emph{American Journal of Clinical Nutrition}, 80:604--613, 2004.

\bibitem{foodcode2026}
U.S. Food and Drug Administration.
\newblock \emph{Food Code 2026}.
\newblock U.S. Department of Health and Human Services, Public Health Service, 2026.

\bibitem{hall2019}
K.~D. Hall et al.
\newblock Ultra-processed diets cause excess calorie intake and weight gain: An inpatient randomized controlled trial of ad libitum food intake.
\newblock \emph{Cell Metabolism}, 30(1):67--77.e3, 2019.
\newblock doi:10.1016/j.cmet.2019.05.008.

\bibitem{jaluria2018}
Y.~Jaluria.
\newblock Manufacture of optical fibers: Drawing and coating processes.
\newblock In \emph{Advanced Materials Processing and Manufacturing}, Mechanical Engineering Series, pages 239--286. Springer, 2018.
\newblock doi:10.1007/978-3-319-76983-7\_8.

\bibitem{krintiras2015}
G.~A. Krintiras, J.~Göbel, A.~J. van~der Goot, and G.~D. Stefanidis.
\newblock Production of structured soy-based meat analogues using simple shear and heat in a Couette Cell.
\newblock \emph{Journal of Food Engineering}, 160:34--41, 2015.
\newblock doi:10.1016/j.jfoodeng.2015.02.015.

\bibitem{kuraishi2001}
C.~Kuraishi, K.~Yamazaki, and Y.~Susa.
\newblock Transglutaminase: Its utilization in the food industry.
\newblock \emph{Food Reviews International}, 17(2):221--246, 2001.
\newblock doi:10.1081/FRI-100001258.

\bibitem{laing1989}
D.~G. Laing and G.~W. Francis.
\newblock The capacity of humans to identify odors in mixtures.
\newblock \emph{Physiology \& Behavior}, 46(5):809--814, 1989.
\newblock doi:10.1016/0031-9384(89)90041-3.

\bibitem{lee2012}
K.~Y. Lee and D.~J. Mooney.
\newblock Alginate: Properties and biomedical applications.
\newblock \emph{Progress in Polymer Science}, 37(1):106--126, 2012.

\bibitem{mason1995}
T.~G. Mason, J.~Bibette, and D.~A. Weitz.
\newblock Elasticity of compressed emulsions.
\newblock \emph{Physical Review Letters}, 75(10):2051--2054, 1995.

\bibitem{mcclements2015}
D.~J. McClements.
\newblock \emph{Food Emulsions: Principles, Practices, and Techniques}, 3rd edition.
\newblock CRC Press, 2015.

\bibitem{mottram1998}
D.~S. Mottram.
\newblock Flavour formation in meat and meat products: A review.
\newblock \emph{Food Chemistry}, 62(4):415--424, 1998.

\bibitem{peters2011}
S.~T. Peters, editor.
\newblock \emph{Composite Filament Winding}.
\newblock ASM International, 2011.

\bibitem{purslow2005}
P.~P. Purslow.
\newblock Intramuscular connective tissue and its role in meat quality.
\newblock \emph{Meat Science}, 70(3):435--447, 2005.

\bibitem{rozin1982}
P.~Rozin.
\newblock ``Taste--smell confusion'' and the duality of the olfactory sense.
\newblock \emph{Perception \& Psychophysics}, 31(4):397--401, 1982.

\bibitem{schutyser2015}
M.~A.~I. Schutyser, P.~J.~M. Pelgrom, A.~J. van~der Goot, and R.~M. Boom.
\newblock Dry fractionation for sustainable production of functional legume protein concentrates.
\newblock \emph{Trends in Food Science \& Technology}, 45(2):327--335, 2015.

\bibitem{singh1996}
N.~Singh and S.~R. Eckhoff.
\newblock Wet milling of corn: A review of laboratory-scale and pilot plant-scale procedures.
\newblock \emph{Cereal Chemistry}, 73(6):659--667, 1996.

\bibitem{vandersman2023}
R.~G.~M. van~der Sman and A.~J. van~der Goot.
\newblock Hypotheses concerning structuring of extruded meat analogs.
\newblock \emph{Current Research in Food Science}, 6:100510, 2023.
\newblock doi:10.1016/j.crfs.2023.100510.

\bibitem{szczesniak1963}
A.~S. Szczesniak, M.~A. Brandt, and H.~H. Friedman.
\newblock Development of standard rating scales for mechanical parameters of texture and correlation between the objective and the sensory methods of texture evaluation.
\newblock \emph{Journal of Food Science}, 28(4):397--403, 1963.
\newblock doi:10.1111/j.1365-2621.1963.tb00217.x.

\bibitem{tetrapak}
Tetra Pak.
\newblock \emph{Dairy Processing Handbook}.
\newblock Tetra Pak Processing Systems AB, online edition, consulted September 2026. Chapters ``Centrifugal separators and milk standardization'' and ``Recombined milk products''. \url{https://dairyprocessinghandbook.tetrapak.com}.

\bibitem{thomasdanguin2014}
T.~Thomas-Danguin, C.~Sinding, S.~Romagny, F.~El~Mountassir, B.~Atanasova, E.~Le~Berre, A.-M. Le~Bon, and G.~Coureaud.
\newblock The perception of odor objects in everyday life: A review on the processing of odor mixtures.
\newblock \emph{Frontiers in Psychology}, 5:504, 2014.
\newblock doi:10.3389/fpsyg.2014.00504.

\bibitem{tornberg2005}
E.~Tornberg.
\newblock Effects of heat on meat proteins: Implications on structure and quality of meat products.
\newblock \emph{Meat Science}, 70(3):493--508, 2005.

\bibitem{vitali2022}
F.~Vitali et al.
\newblock Semantics of dairy fermented foods: A microbiologist's perspective.
\newblock \emph{Foods}, 11:1939, 2022.

\bibitem{yang2027}
S.~Yang, Y.~Xu, C.~Jia, and Y.~Wang.
\newblock Toward trustworthy LLM-based scientific information extraction: An ontology-guided, validation-aware workflow and a single-domain case study.
\newblock \emph{Information Processing and Management}, 64:105179, 2027.
\newblock Available online 22 September 2026. doi:10.1016/j.ipm.2026.105179.

\bibitem{zink2024}
J.~I. Zink, V.~Lutz-Bueno, S.~Handschin, C.~Dütsch, A.~Diaz, P.~Fischer, and E.~J. Windhab.
\newblock Structural and mechanical anisotropy in plant-based meat analogues.
\newblock \emph{Food Research International}, 179:113968, 2024.
\newblock doi:10.1016/j.foodres.2024.113968.

\end{thebibliography}


\end{document}
